\documentclass[12pt,a4paper]{article}

\usepackage[utf8]{inputenc}
\usepackage[T5]{fontenc}
\usepackage[utf8]{vntex}

\usepackage{lmodern}
\usepackage{geometry}
\geometry{margin=2.5cm}
\usepackage{graphicx}
\graphicspath{{figures/}{./}}
\usepackage{amsmath}
\usepackage{booktabs}
\usepackage{float}
\usepackage{array}
\setlength{\parindent}{0pt}
\setlength{\parskip}{6pt}
\usepackage{caption}
\usepackage{subcaption}
\usepackage{xcolor}
\usepackage[hypertexnames=false]{hyperref}
\usepackage{longtable}

% Ten muc caption tieng Viet (vntex cai dat sau begin-document, nen ta cung
% dat trong cung mot hook de ghi de).
\AtBeginDocument{%
  \renewcommand{\contentsname}{Mục lục}%
  \renewcommand{\listfigurename}{Danh sách hình vẽ}%
  \renewcommand{\listtablename}{Danh sách bảng}%
  \renewcommand{\figurename}{Hình}%
  \renewcommand{\tablename}{Bảng}%
  \renewcommand{\refname}{Tài liệu tham khảo}%
}

% Giu float trong pham vi muc cua no
\usepackage[section]{placeins}

\hypersetup{
    colorlinks=true,
    linkcolor=black,
    urlcolor=blue,
}

\setlength{\floatsep}{10pt plus 2pt minus 2pt}
\setlength{\textfloatsep}{12pt plus 2pt minus 4pt}
\setlength{\intextsep}{10pt plus 2pt minus 2pt}
\renewcommand{\textfraction}{0.05}
\renewcommand{\floatpagefraction}{0.85}
\setlength{\emergencystretch}{3em}
\hbadness=2000

% Ho tro hinh: chen anh that khi co file, neu khong thi chen mot khung giu cho,
% de bao cao van bien dich duoc truoc khi co so do. Khi tha file .png dung ten
% vao thu muc nay, hinh se tu dong hien len.
\newcommand{\reportfig}[3]{%
  \IfFileExists{figures/#2}{%
    \includegraphics[width=#1\textwidth]{#2}%
  }{%
    \fbox{\begin{minipage}[c][5.5cm][c]{#1\textwidth}%
      \centering\itshape
      [\,Hình sẽ được chèn\,]\par\vspace{6pt}
      #3\par\vspace{10pt}%
      \normalfont\ttfamily\footnotesize #2%
    \end{minipage}}%
  }%
}

% Ho tro logo: y tuong tuong tu cho logo trang bia, de mot file logo thieu
% khong lam hong qua trinh bien dich.
\newcommand{\titlelogo}[1]{%
  \IfFileExists{#1}{\includegraphics[width=0.9\textwidth]{#1}}{%
    \fbox{\parbox[c][2.6cm][c]{0.8\textwidth}{\centering\ttfamily\footnotesize #1}}%
  }%
}

\begin{document}
\raggedbottom
\pagenumbering{gobble}

% =========================================================================
% TRANG BIA
% =========================================================================
\begin{titlepage}
    \centering
    \textbf{\large Trường Đại học Công nghệ, ĐHQGHN, Hà Nội}\\[0.2cm]
    \textbf{\large Viện Trí tuệ Nhân tạo}\\[0.5cm]
    \rule{10cm}{0.4pt}\\

    \vspace{2cm}
    \begin{minipage}{0.45\textwidth}
        \begin{center}
            \titlelogo{uet_logo.jpg}
        \end{center}
    \end{minipage}
    \hfill
    \begin{minipage}{0.45\textwidth}
        \begin{center}
            \titlelogo{ai_logo.png}
        \end{center}
    \end{minipage}

    \vspace{1.5cm}
    {\large \textbf{BÁO CÁO MÔN HỌC}}\\[0.5cm]
    {\large \textbf{THỰC TẬP DOANH NGHIỆP}}\\[0.5cm]

    \vspace{5pt}
    {\large{
    Đề tài: \textbf{HumJob: Bộ công cụ trên trình duyệt, ưu tiên xử lý cục bộ,
    cho giọng hát - Ký âm, Phân tích và Luyện tai/Luyện giọng}
    }}\\[0.25cm]

    \vspace{12pt}
    {\large
    \begin{tabular}{>{\bfseries}l l}
        Giảng viên hướng dẫn: & TS. Trần Hồng Việt \\
    \end{tabular}
    }\\[0.5cm]

    \vspace{12pt}
    {\large
    \begin{tabular}{>{\bfseries}l l}
        Sinh viên thực hiện: & Vũ Nguyên Đan - 23020351 \\
    \end{tabular}
    }\\[0.5cm]

    \vfill
    \centering
    \textbf{Hà Nội, 2026}
\end{titlepage}

% =========================================================================
% TRANG NHAN XET
% =========================================================================
\newpage
\begin{center}
    \textbf{Nhận xét của giảng viên hướng dẫn:}
\end{center}

\vspace{0.2cm}

\noindent\makebox[\linewidth]{\dotfill}\\[0.2cm]
\noindent\makebox[\linewidth]{\dotfill}\\[0.2cm]
\noindent\makebox[\linewidth]{\dotfill}\\[0.2cm]
\noindent\makebox[\linewidth]{\dotfill}\\[0.2cm]
\noindent\makebox[\linewidth]{\dotfill}\\[0.2cm]
\noindent\makebox[\linewidth]{\dotfill}\\[0.2cm]
\noindent\makebox[\linewidth]{\dotfill}\\[0.2cm]
\noindent\makebox[\linewidth]{\dotfill}\\[0.2cm]
\noindent\makebox[\linewidth]{\dotfill}\\[0.2cm]
\noindent\makebox[\linewidth]{\dotfill}\\[0.2cm]
\noindent\makebox[\linewidth]{\dotfill}\\[0.2cm]
\noindent\makebox[\linewidth]{\dotfill}\\[0.2cm]
\noindent\makebox[\linewidth]{\dotfill}\\[0.2cm]
\noindent\makebox[\linewidth]{\dotfill}\\[0.2cm]
\noindent\makebox[\linewidth]{\dotfill}\\[0.2cm]
\noindent\makebox[\linewidth]{\dotfill}\\[0.2cm]
\noindent\makebox[\linewidth]{\dotfill}\\[0.2cm]
\noindent\makebox[\linewidth]{\dotfill}

\noindent\hspace{3cm} Điểm (bằng số) \dotfill{} Điểm (bằng chữ) \dotfill

\vspace{1.5cm}

\begin{center}
    Hà Nội, \ldots{} \ldots{} 2026\\[0.2cm]
    \textbf{Giảng viên hướng dẫn}\\[0.1cm]
    (Ký và ghi rõ họ tên)
\end{center}

% =========================================================================
% LOI CAM ON
% =========================================================================
\newpage
\begin{center}
    \textbf{\large LỜI CẢM ƠN}
\end{center}

Em xin cảm ơn giảng viên hướng dẫn vì những chỉ dẫn và góp ý trong suốt môn học
Thực tập Doanh nghiệp, những điều đã định hình cả hướng đi của công trình này,
\textbf{HumJob}, lẫn chuẩn mực về sự trung thực học thuật mà em cố gắng giữ cho nó.

Báo cáo này là công trình của một sinh viên duy nhất. Với phạm vi rộng của hệ
thống và thời gian có hạn, chắc chắn nó còn nhiều thiếu sót. Đặc biệt, kết luận
trung tâm của phần đánh giá bộ ký âm là một điều đáng suy ngẫm: bộ chỉ số tổng
hợp trên dữ liệu tổng hợp mà bộ ký âm được tinh chỉnh theo đã bão hòa ở mức trần,
trong khi hệ thống chưa được kiểm chứng trên đầu vào ngân nga thực. Em chọn báo
cáo điều đó một cách trung thực, và xây dựng lại phần đánh giá dựa trên các bằng
chứng phong phú hơn, thay vì trình bày một chỉ số đã bão hòa như một thành công.
Em rất mong nhận được nhận xét và góp ý của giảng viên để công trình có thể hoàn
thiện hơn.

% =========================================================================
% TOM TAT
% =========================================================================
\newpage
\begin{center}
    \textbf{\large TÓM TẮT}
\end{center}

\textbf{HumJob} là một bộ công cụ chạy trên trình duyệt, ưu tiên xử lý cục bộ,
dành cho giọng hát. Nó tập hợp sáu công cụ cùng chia sẻ một nền tảng, đó là phát
hiện cao độ áp dụng cho giọng người, và tất cả đều chạy trên chính máy của người
dùng: một \emph{bộ ký âm giai điệu} biến một giai điệu ngân nga thành nốt nhạc,
giọng và hợp âm gợi ý, rồi xuất ra dưới dạng MIDI, MusicXML và bản nhạc được khắc;
một \emph{bộ tìm cao độ} báo cáo giọng, nhịp độ và mã Camelot của bất kỳ đoạn âm
thanh nào, kể cả nhạc đa âm; một \emph{bộ dịch giọng} chuyển một bản nhạc sang
giọng mới; một \emph{bộ giám sát cao độ thời gian thực} và luyện giọng; một
\emph{bộ hát theo (sing-along)} chấm điểm phần trình bày so với một giai điệu tham
chiếu; và một \emph{bộ luyện tai} cho nhận diện quãng, hợp âm, thang âm và bậc
thang âm. Một \emph{bảng luyện tập (practice hub)} gộp lịch sử luyện tập mà các
công cụ ghi lại cục bộ thành các xu hướng tiến bộ. Một số tính năng có thêm hỗ trợ
tùy chọn từ mô hình ngôn ngữ, gồm gợi ý luyện tập bằng ngôn ngữ tự nhiên, kế hoạch
luyện tập hằng tuần, chỉnh sửa bản nhạc bằng ngôn ngữ tự nhiên, viết lại toàn bộ
giai điệu, và phối lại hợp âm, đây là phần duy nhất của hệ thống từng rời khỏi máy
người dùng.

Phần lõi kỹ thuật của bộ công cụ, cũng là đối tượng của phần đánh giá chính thức
trong báo cáo này, là bộ ký âm. Nó giải quyết bài toán ký âm giai điệu đơn âm từ
giọng người trần, một bài toán khó vì một đường cao độ liên tục phải được cắt
thành các nốt rời rạc với cao độ, thời điểm khởi phát (onset) và trường độ đúng.
HumJob làm cho việc này khả thi nhờ hai ràng buộc giao diện có chủ đích. Người
dùng ngân nga âm tiết ``da-da-da'', trong đó phụ âm tạo cho mỗi nốt một onset rõ
ràng và tách được các nốt cùng cao độ lặp lại; và người dùng ngân nga theo một máy
đếm nhịp với nhịp độ đã biết, biến nhịp điệu từ bài toán ước lượng thành bài toán
gióng vào một lưới đã biết. Đường ống (pipeline) được xây dựng như một chuỗi các
hàm thuần trên ba kiểu dữ liệu, với mọi tham số điều chỉnh được đặt tập trung, và
hỗ trợ năm backend phát hiện cao độ có thể thay thế cho nhau (pyin, CREPE, PESTO,
FCNF0++ và basic-pitch). Một bộ phân đoạn nhận biết lưới nhịp và một bước hợp nhất
chạy cho mọi backend xử lý dạng lỗi chủ đạo, còn một bước hiệu chỉnh quãng tám dựa
trên phổ sửa một lỗi có hệ thống của bộ theo dõi cao độ trên các câu nhạc legato.

Bộ ký âm được đánh giá trên một khung dữ liệu tổng hợp có nhãn, chấm điểm phần phát
hiện nốt (F1 của cao độ và onset) và, riêng biệt, phần nhịp điệu (độ chính xác của
onset và trường độ đã lượng tử hóa so với một lưới dự kiến). Trên các mẫu ``thực
tế'' giàu biểu cảm, F1 của nốt tăng từ $0.799$ lên $1.000$ qua ba vòng cải tiến, và
khung đo bão hòa; nửa nhịp điệu cũng bão hòa khi một bước tinh chỉnh nhịp độ được
thêm vào. Một chỉ số F1 ở mức trần không thể phân biệt một bộ phân đoạn tốt với một
bộ kém hơn, và nó hoàn toàn là tổng hợp, nên phần đánh giá được xây dựng lại quanh
bốn loại bằng chứng không cần một ca sĩ chính xác, vì tác giả không phải là ca sĩ
và không tồn tại tập dữ liệu ngân nga thực có nhãn. Đó là: các tỷ lệ lỗi chẩn đoán
trên một kho ngữ liệu lớn và đa dạng hơn, giúp phân biệt trở lại và khu trú dạng
lỗi chủ đạo với đầu vào biểu cảm về các nốt cùng cao độ bị gộp (khả năng thu hồi
lặp giảm còn khoảng $0.61$); các tính chất đúng đắn kiểu biến hình (dịch giọng,
khuếch đại, nhiễu, đệm khoảng lặng và bất biến nhịp độ) không cần giai điệu tham
chiếu và vẫn đúng trên cả âm thanh thực; một bước tái tạo khứ hồi (round-trip) chấm
điểm một đoạn ngân nga thực so với chính âm thanh được tổng hợp lại của nó (một tín
hiệu kiểm tra không giám sát, trung bình $0.62$ trên bốn đoạn thu thực); và một bước
hiệu chuẩn độ ``thực tế'' của dữ liệu tổng hợp so với các đoạn thu thực, khẳng định
rằng tần số rung (vibrato) và độ trôi cao độ được giả định là gần với giọng thực.
Một phát hiện hữu ích và trung thực là đường ống này không bị chia nhỏ quá mức
(over-split) dưới nhiễu loạn mạnh như từng dự đoán; các hàng rào phòng vệ của nó
vẫn giữ vững, và thay vào đó nó suy giảm bằng cách gộp hoặc bỏ nốt. Không loại nào
trong bốn loại là một con số về độ chính xác thực tế, điều đó vẫn cần nhãn chuẩn
thực, nhưng cùng nhau chúng thay thế một con số đã bão hòa bằng một bức tranh vẫn
còn hữu ích. Báo cáo này ghi lại toàn bộ bộ công cụ, tập trung phần đánh giá vào bộ
ký âm, và trình bày trung thực điều mà mỗi loại bằng chứng có thể và không thể
chứng minh.

% =========================================================================
% MUC LUC
% =========================================================================
\newpage
\pagenumbering{arabic}
\tableofcontents
\newpage
\listoffigures
\newpage
\listoftables
\newpage

% =========================================================================
% MUC 1: GIOI THIEU
% =========================================================================
\section{Giới thiệu}

\subsection{Động lực và ý nghĩa}

Có một tình huống âm nhạc quen thuộc mà không có công cụ tốt nào giải quyết. Một
người có một giai điệu trong đầu, ngân nga nó ra, và muốn nó được ghi lại: dưới
dạng nốt trên khuông nhạc, một tệp MIDI để chỉnh sửa, hay một phác thảo hợp âm để
chơi lại. Việc biến đoạn ngân nga đó thành bản nhạc chính là bài toán
\emph{ký âm theo lời ngân nga (query-by-humming)}, và nó khó hơn vẻ ngoài. Giọng
người tạo ra một đường cao độ liên tục duy nhất, trượt, dao động và trôi; ký âm đòi
hỏi đường cong đó phải được cắt thành một chuỗi nốt rời rạc, mỗi nốt có cao độ xác
định, thời điểm bắt đầu xác định và một trường độ ghi được, đồng thời cần một giọng
và một lưới nhịp mà trong đó các trường độ ấy có ý nghĩa.

Hai họ công cụ hiện có né tránh một phần bài toán thay vì giải nó. Các bộ chuyển âm
thanh sang MIDI tổng quát được huấn luyện trên nhạc cụ và nhạc đa âm; trên giọng
người trần, chúng dễ nhảy quãng tám và không tận dụng bất kỳ cấu trúc nào mà một
đoạn ngân nga có chủ đích cung cấp. Các bộ giám sát cao độ và bộ lên dây thời gian
thực báo cáo cao độ tức thời nhưng không bao giờ đưa ra một phép phân đoạn thành
nốt. Điều còn thiếu là một hệ thống được xây dựng riêng cho một con người hợp tác,
sẵn lòng ngân nga theo cách giúp việc ký âm trở nên khả thi.

HumJob chính là hệ thống đó. Quyết định thiết kế trung tâm của nó là chấp nhận hai
ràng buộc nhỏ về cách người dùng ngân nga, và xây dựng toàn bộ đường ống quanh
chúng:

\begin{itemize}
    \item Người dùng ngân nga âm tiết \textbf{``da-da-da''}. Phụ âm khép nhẹ thanh
    quản trong chốc lát ở đầu mỗi nốt, tạo ra một onset biên độ rõ nét và, quan
    trọng hơn, tách hai nốt cùng cao độ lặp lại mà một đoạn ngân nga legato sẽ nhòe
    thành một.
    \item Người dùng ngân nga \textbf{theo một máy đếm nhịp với nhịp độ đã biết}.
    Một lưới nhịp đã biết biến nhịp điệu từ một bài toán ước lượng khó thành một bài
    toán gióng: onset và trường độ của nốt được gióng vào vị trí lưới gần nhất thay
    vì đoán từ thời gian thô.
\end{itemize}

Hai ràng buộc này là cột chống đỡ. Chúng là thứ khiến bài toán trở nên khả thi cho
một đồ án một sinh viên, và một chủ đề lặp lại trong báo cáo này là các lỗi còn lại
của bộ ký âm là lỗi ở việc chưa khai thác hết chúng, chứ không phải lỗi của các mô
hình cao độ nền tảng.

\subsection{Từ một bộ ký âm đến một bộ công cụ cho giọng hát}
\label{sec:toolkit-intro}

HumJob khởi đầu là bộ ký âm đó, và nó vẫn là lõi kỹ thuật cũng như trọng tâm đánh
giá của báo cáo này. Trong quá trình xây dựng, một họ công cụ liên quan chặt chẽ
dành cho giọng hát đã nảy sinh từ cùng một nền tảng, bởi mỗi công cụ đều dựa trên
cùng một thao tác: ước lượng cao độ của giọng người, cục bộ và ngay trong trình
duyệt. Hệ thống hoàn chỉnh vì thế là một bộ công cụ gồm sáu công cụ, trình bày dưới
dạng các thẻ (tab) trong một ứng dụng web ưu tiên xử lý cục bộ:

\begin{itemize}
    \item một \textbf{Bộ ký âm (Transcriber)} biến đoạn ngân nga thành nốt, giọng
    và hợp âm, đây là đường ống lõi, kèm một trình soạn thảo khuông nhạc tương tác
    để sửa bản nháp và phối lại hợp âm bằng phương pháp cục bộ hoặc bằng mô hình
    ngôn ngữ;
    \item một \textbf{Bộ tìm cao độ (Pitch Finder)} báo cáo giọng, nhịp độ và mã
    Camelot của một đoạn tải lên, kể cả nhạc đa âm;
    \item một \textbf{Bộ dịch giọng (Transposer)} chuyển một bản nhạc, từ tệp hoặc
    từ đoạn ngân nga gần nhất, sang giọng mới;
    \item một \textbf{Bộ thời gian thực (Realtime)} giám sát và luyện giọng: hiển
    thị cao độ trực tiếp, một bộ lên dây đàn guitar và một tập bài luyện tập;
    \item một \textbf{Bộ hát theo (Sing-Along)} phát một giai điệu tham chiếu và
    chấm điểm phần trình bày từng nốt, kèm gợi ý luyện tập tùy chọn bằng mô hình
    ngôn ngữ;
    \item một \textbf{Bộ luyện tai (Ear Trainer)} phát quãng, hợp âm, thang âm và
    tiến trình kết (cadence) để người dùng nhận diện.
\end{itemize}

Một thẻ thứ bảy, \textbf{bảng luyện tập (practice hub)}, là một bảng điều khiển
trang chủ hơn là một công cụ theo nghĩa trên: nó gộp lịch sử luyện tập mà các thẻ
khác ghi lại cục bộ thành các xu hướng tiến bộ ngoại tuyến. Một số thẻ cũng cung
cấp hỗ trợ tùy chọn từ mô hình ngôn ngữ, mô tả ở Mục~\ref{sec:coaching}.

Báo cáo này mô tả toàn bộ bộ công cụ nhưng tập trung phần đánh giá vào bộ ký âm, vì
đó là nơi có đóng góp nghiên cứu và bài toán khó nhất, đo được rõ nhất của đồ án.
Năm công cụ còn lại là các sản phẩm kỹ thuật xây trên nền tảng chung;
Mục~\ref{sec:toolkit} mô tả chúng và Phụ lục~A minh họa từng cái.

\subsection{Mục tiêu}

Đồ án theo đuổi các mục tiêu sau:

\begin{itemize}
    \item Thiết kế và hiện thực một đường ống ký âm giai điệu hoàn chỉnh, ưu tiên
    xử lý cục bộ, nhận một bản thu ngân nga và trả về nốt, giọng và hợp âm gợi ý,
    xuất ra MIDI, MusicXML và bản nhạc được khắc.
    \item Cấu trúc đường ống thành một chuỗi các hàm thuần trên một mô hình dữ liệu
    nhỏ và cố định, để mỗi bước kiểm thử được độc lập và mọi tham số điều chỉnh
    được đặt tập trung thay vì rải rác dưới dạng các con số ma thuật.
    \item Hỗ trợ nhiều backend phát hiện cao độ có thể thay thế nhau sau một giao
    diện duy nhất, để mô hình sinh nốt có thể được chọn hoặc thay mà không đụng đến
    phần còn lại của đường ống.
    \item Cung cấp đường ống dưới cả dạng công cụ dòng lệnh lẫn ứng dụng web ghi âm
    cục bộ và gọi cùng một đoạn mã, kèm một trình soạn thảo tương tác để sửa bản
    nháp tự động.
    \item Tái sử dụng cùng nền tảng phát hiện cao độ ưu tiên cục bộ trên một họ
    công cụ liên quan cho giọng hát, gồm phân tích, dịch giọng, giám sát thời gian
    thực, luyện giọng và luyện tai, để lõi phục vụ cho nhiều hơn một nhiệm vụ âm
    nhạc.
    \item Xây dựng một khung đánh giá đo bộ ký âm một cách trung thực, và nêu rõ
    khung đó có thể và không thể chứng minh điều gì.
\end{itemize}

\subsection{Câu hỏi nghiên cứu}

Đồ án được tổ chức quanh ba câu hỏi nghiên cứu. Chúng liên quan riêng đến bộ ký âm,
vì đó là nơi có công việc khảo cứu; các công cụ khác được đánh giá như những sản
phẩm chức năng (Mục~\ref{sec:toolkit}, Phụ lục~A).

\begin{itemize}
    \item \textbf{RQ1 - Các ràng buộc giao diện có giúp ích không?} Hai ràng buộc
    có chủ đích (một onset phụ âm mỗi nốt, và ngân nga theo một lưới nhịp độ đã
    biết) có biến ký âm giai điệu từ một bài toán ước lượng thành một bài toán phân
    đoạn và gióng khả thi hay không?
    \item \textbf{RQ2 - Đây là bài toán phân đoạn hay bài toán mô hình?} Khi nốt
    ra sai, lỗi chủ yếu nằm ở phân đoạn và lượng tử hóa của một đường cao độ vốn
    đúng, hay ở chính mô hình cao độ?
    \item \textbf{RQ3 - Hệ thống có thể được đánh giá tốt tới đâu?} Hệ thống khớp
    với nhãn chuẩn tốt tới đâu, và giới hạn phương pháp luận của phép đo đó là gì,
    khi đồ án chỉ có nhãn chuẩn tổng hợp và không có kho ngữ liệu ngân nga thực có
    nhãn?
\end{itemize}

% =========================================================================
% MUC 2: CO SO LY THUYET
% =========================================================================
\section{Cơ sở lý thuyết}

\subsection{Ký âm giai điệu và bài toán con đơn âm}

Ký âm nhạc tự động nói chung là việc chuyển một tín hiệu âm thanh thành một bản
nhạc ký hiệu. Bài toán tổng quát là đa âm và rất khó. HumJob giải quyết trường hợp
hạn chế \emph{đơn âm}: một giọng, một nốt tại một thời điểm. Ràng buộc này loại bỏ
nhu cầu tách các cao độ đồng thời, nhưng không làm bài toán trở nên tầm thường, vì
khó khăn dịch từ tách nguồn sang \emph{phân đoạn}: quyết định nốt này kết thúc và
nốt kia bắt đầu ở đâu, và mỗi nốt dài bao nhiêu theo nghĩa nhịp.

Một đường ống ký âm đơn âm có ba giai đoạn logic: ước lượng một đường cao độ theo
thời gian, cắt đường cong đó thành các nốt rời rạc, và đặt các nốt đó lên một giọng
và một lưới nhịp. HumJob theo đúng hình dạng này, và phần lớn công sức kỹ thuật nằm
ở giai đoạn thứ hai và thứ ba.

\subsection{Ước lượng tần số cơ bản}

Giai đoạn đầu ước lượng tần số cơ bản ($f_0$) của giọng, theo từng khung. Có vài họ
bộ ước lượng và HumJob hỗ trợ năm trong số đó sau một giao diện:

\begin{itemize}
    \item \textbf{pYIN}~\cite{pyin} là một bộ ước lượng xử lý tín hiệu cổ điển, bổ
    sung cho thuật toán YIN các ngưỡng xác suất và một bước giải mã Viterbi. Nó
    nhanh, nhẹ về phụ thuộc, và là mặc định của đồ án.
    \item \textbf{CREPE}~\cite{crepe} là một mạng nơ-ron tích chập được huấn luyện
    để ước lượng cao độ trực tiếp từ dạng sóng, và ổn định trên giọng hát.
    \item \textbf{PESTO}~\cite{pesto} là một bộ ước lượng cao độ tự giám sát, đẳng
    biến theo phép dịch giọng; nó là mặc định của ứng dụng web vì chính xác trên
    giọng trong khi nhẹ và nhanh hơn CREPE.
    \item \textbf{FCNF0++}~\cite{penn} là một bộ ước lượng cao độ và tính chu kỳ
    hoàn toàn tích chập, ngang tầm độ chính xác với PESTO.
    \item \textbf{basic-pitch}~\cite{basicpitch} là một mô hình ký âm nốt bằng
    nơ-ron được huấn luyện trên nhạc cụ, xuất trực tiếp các sự kiện nốt, bỏ qua
    đường cong cấp khung.
\end{itemize}

Sự thật nền quan trọng cho báo cáo này liên quan đến một lỗi cụ thể của các bộ theo
dõi giải mã bằng Viterbi. Trên một câu nhạc legato liên tục có tiếng, không có
khoảng lặng để đặt lại bộ giải mã, tiên nghiệm chuyển tiếp ``đứng yên'' của bộ giải
mã Viterbi có thể khóa cả một nốt vào sóng hài phụ của nó, báo cáo nó thấp đúng một
quãng tám. Đây không phải là hiện tượng nhiễu âm thanh; nó là hiện tượng do giải mã,
và Mục~\ref{sec:octave} mô tả cách HumJob phát hiện và sửa nó.

\subsection{Từ đường cao độ đến nốt}

Biến một đường cong thành nốt đòi hỏi ba quyết định con. \emph{Xác định tiếng
(voicing)} quyết định khung nào thực sự mang cao độ, loại bỏ tiếng thở và khoảng
lặng. \emph{Phân đoạn} quyết định ranh giới giữa các nốt nằm ở đâu, mà với một đoạn
ngân nga ``da-da-da'' nghĩa là tìm những chỗ trũng năng lượng ngắn do phụ âm tạo
ra. \emph{Hợp nhất} làm điều ngược lại khi cần: một nốt giữ hơi duy nhất bị vibrato
xé thành nhiều mảnh phải được nối lại thành một. Phân đoạn và hợp nhất kéo về hai
hướng đối nghịch, và cân bằng chúng là khó khăn cốt lõi của đường ống.

\subsection{Tìm giọng và lượng tử hóa nhịp điệu}

Khi đã có nốt, còn hai quyết định âm nhạc. \emph{Tìm giọng} xác định thang âm mà
giai điệu nằm trong; HumJob tương quan phân bố lớp cao độ đã hát với các hồ sơ âm
thử (probe-tone) Krumhansl-Kessler cho cả hai mươi tư giọng trưởng và
thứ~\cite{krumhansl}. \emph{Lượng tử hóa} gióng onset và trường độ của nốt vào một
lưới nhịp suy ra từ nhịp độ đã biết, đây là điều cho phép trường độ được viết thành
các giá trị nốt gọn thay vì các phân số vụng về. Khi nhịp độ sai, lượng tử hóa tạo
ra các trường độ không nguyên mà phần mềm ký âm hiển thị thành chuỗi nốt nối
(tie), đó là triệu chứng nhìn thấy được của hầu hết các lỗi nhịp điệu.

\subsection{Tại sao máy đếm nhịp và ``da-da-da''}

Hai ràng buộc giới thiệu ở Mục~1 là lời đáp của đồ án cho hai bài toán con khó nhất
ở trên. Phụ âm trong ``da-da-da'' cung cấp trực tiếp gợi ý onset mà phân đoạn cần,
và đặc biệt nó là cách đáng tin duy nhất để tách hai nốt cùng cao độ liên tiếp. Máy
đếm nhịp cung cấp lưới mà lượng tử hóa gióng vào. Không ràng buộc nào là gánh nặng
cho một người dùng hợp tác, và cùng nhau chúng thay hai bài toán ước lượng bằng hai
bài toán phát hiện. Đồ án xem chúng là các giả định cố định, không phải là các hạn
chế cần tối ưu để loại bỏ.

% =========================================================================
% MUC 3: THIET KE HE THONG VA PHUONG PHAP
% =========================================================================
\section{Thiết kế hệ thống và phương pháp}

\subsection{Tổng quan kiến trúc}

Bộ ký âm được tổ chức thành một chuỗi các hàm thuần trên ba kiểu dữ liệu.
Hình~\ref{fig:arch} cho thấy cấu trúc cấp cao: một gói Python lõi hiện thực đường
ống, một công cụ dòng lệnh gọi nó trực tiếp, và một ứng dụng web FastAPI ghi âm
trong trình duyệt rồi gọi cùng đường ống ở phía máy chủ. Cùng ứng dụng web này cũng
chứa năm công cụ còn lại của bộ công cụ dưới dạng các thẻ trình duyệt riêng
(Mục~\ref{sec:toolkit}); mục này mô tả đường ống ký âm, vốn là lõi kỹ thuật chung.
Mọi thứ chạy cục bộ, và không có âm thanh nào từng được gửi tới bất kỳ dịch vụ bên
thứ ba nào. Ngoại lệ với nguyên tắc chỉ chạy cục bộ là một họ nhỏ các tính năng AI
tùy chọn (Mục~\ref{sec:coaching}): chỉ khi người dùng chủ động nhấn một nút thì một
trong số chúng mới gửi một bản tóm tắt văn bản hoặc số ngắn, không bao giờ là âm
thanh, tới một API mô hình ngôn ngữ bên ngoài.

\begin{figure}[H]
\centering
\reportfig{0.9}{architecture.png}{Kiến trúc hệ thống: gói đường ống lõi, điểm vào
dòng lệnh, và ứng dụng web FastAPI cùng giao diện trình duyệt tĩnh. Cùng đoạn mã
đường ống phục vụ cả hai điểm vào.}
\caption{Kiến trúc hệ thống cấp cao của HumJob.}
\label{fig:arch}
\end{figure}

Ba kiểu dữ liệu là hợp đồng mà mọi bước chia sẻ (Bảng~\ref{tab:types}). Một
\texttt{Frame} là một khung phân tích dày, khoảng một khung mỗi bước nhảy $11.6$ ms,
mang thời điểm khung, $f_0$ ước lượng, một độ tin cậy về tiếng, và một năng lượng
ngắn hạn. Một \texttt{NoteEvent} là một nốt rời rạc sau phân đoạn, mang thời điểm
bắt đầu và kết thúc theo giây, cao độ nguyên, cao độ phân số thô, một độ lệch cent
so với nửa cung gần nhất, và, khi đã lượng tử hóa, vị trí và trường độ được gióng
lưới theo đơn vị nốt đen. Một \texttt{Score} là toàn bộ bản ký âm, sẵn sàng để xuất.

\begin{table}[H]
\centering
\caption{Ba kiểu dữ liệu mà đường ống được xây dựng trên đó.}
\label{tab:types}
\begin{tabular}{l p{0.60\textwidth}}
\toprule
\textbf{Kiểu} & \textbf{Vai trò} \\
\midrule
\texttt{Frame}     & Một khung phân tích ($\approx 11.6$ ms): thời điểm khung, $f_0$
(NaN nếu không có tiếng), độ tin cậy về tiếng, năng lượng ngắn hạn. Là thứ bộ theo
dõi cao độ nhìn thấy. \\
\texttt{NoteEvent} & Một nốt rời rạc sau phân đoạn: thời điểm bắt đầu/kết thúc theo
giây, cao độ MIDI nguyên, cao độ phân số thô, độ lệch cent, và (sau lượng tử hóa)
vị trí lưới và trường độ theo đơn vị nốt đen. \\
\texttt{Score}     & Toàn bộ bản ký âm: chuỗi nốt, giọng phát hiện được, độ lệch
chỉnh dây, và hợp âm, sẵn sàng để xuất. \\
\bottomrule
\end{tabular}
\end{table}

Một quy tắc thiết kế mà đồ án giữ vững là mọi tham số điều chỉnh đều nằm trong một
đối tượng tham số tập trung duy nhất, không phải là một con số ma thuật bên trong
một bước. Thay đổi hành vi nghĩa là đổi một núm có tên, điều giữ cho các bước dễ đọc
và các thí nghiệm tái lập được.

\subsection{Đường ống ký âm}

Hình~\ref{fig:pipeline} cho thấy các bước của đường ống theo thứ tự. Chuỗi là:

\begin{center}
âm thanh $\rightarrow$ tiền xử lý $\rightarrow$ [sinh nốt] $\rightarrow$ hợp nhất
$\rightarrow$ chỉnh dây $\rightarrow$ giọng $\rightarrow$ lượng tử hóa $\rightarrow$
hợp âm $\rightarrow$ xuất.
\end{center}

\begin{figure}[H]
\centering
\reportfig{0.92}{pipeline-transcriber.png}{Đường ống ký âm. Sinh nốt có hai nhánh:
các bộ theo dõi cấp khung (pyin / CREPE / PESTO / FCNF0++) đưa vào hiệu chỉnh quãng
tám, xác định tiếng và bộ phân đoạn nhận biết lưới, còn basic-pitch xuất trực tiếp
các sự kiện nốt. Cả hai nhánh sau đó dùng chung hợp nhất, chỉnh dây, tìm giọng,
lượng tử hóa, hợp âm và xuất.}
\caption{Đường ống ký âm của HumJob, cả hai nhánh sinh nốt.}
\label{fig:pipeline}
\end{figure}

\textbf{Sinh nốt} có hai nhánh, chọn theo một tham số backend. Backend nơ-ron xuất
sự kiện nốt (basic-pitch) xuất nốt trực tiếp và không có đường cong khung. Các bộ
theo dõi cấp khung (pyin, CREPE, PESTO, FCNF0++) tạo ra một đường cong \texttt{Frame}
rồi đi qua hiệu chỉnh quãng tám, một cổng xác định tiếng, và bộ phân đoạn nhận biết
lưới để thành nốt.

\textbf{Hợp nhất} chạy cho mọi backend. Nó nối lại các mảnh mà một nốt giữ hơi, đầy
vibrato để lại, đây là cách sửa lỗi ``một nốt biến thành nhiều mảnh vụn''. Nó nhận
biết lưới: nó sẽ không nối hai nốt cùng cao độ băng qua một onset lưới, nên nó không
bao giờ hoàn tác một sự tái phát âm mà bộ phân đoạn đã cố ý tách trên phách.

\textbf{Các bước còn lại} áp một hiệu chỉnh chỉnh dây toàn cục, tìm giọng bằng tương
quan Krumhansl-Kessler, lượng tử hóa onset và trường độ vào lưới nhịp độ, gợi ý hợp
âm, và xuất ra MIDI, MusicXML và bản nhạc được khắc.

\subsection{Các backend phát hiện nốt}

Năm backend đánh đổi độ chính xác, tốc độ và miền huấn luyện theo những cách khác
nhau. Bảng~\ref{tab:backends} tóm tắt chúng. Mặc định khác nhau theo điểm vào, và
điều đó là có chủ đích: mặc định của thư viện là pyin (nhanh và nhẹ phụ thuộc), mặc
định của dòng lệnh là basic-pitch, và mặc định của ứng dụng web là PESTO (chính xác
nhất trên giọng). Bốn bộ theo dõi cấp khung đều đưa vào cùng một bộ phân đoạn; chỉ
basic-pitch bỏ qua nó.

\begin{table}[H]
\centering
\caption{Các backend phát hiện nốt.}
\label{tab:backends}
\begin{tabular}{l p{0.34\textwidth} p{0.34\textwidth}}
\toprule
\textbf{Backend} & \textbf{Là gì} & \textbf{Hợp nhất cho / ghi chú} \\
\midrule
\texttt{pesto} & Bộ ước lượng cao độ tự giám sát, đẳng biến theo phép dịch giọng &
Chính xác nhất trên giọng; mặc định ứng dụng web; nhẹ và nhanh hơn CREPE \\
\texttt{fcnf0} & FCNF0++, một bộ ước lượng $f_0$ hoàn toàn tích chập & Ngang tầm độ
chính xác với PESTO; nổi bật trên giọng thực; cài đặt tùy chọn \\
\texttt{crepe} & Mạng nơ-ron tích chập cao độ, huấn luyện trên giọng & Ổn định trên
ngân nga và hát; dùng mô hình full chính xác với giải mã Viterbi \\
\texttt{basic\_pitch} & Mạng tích chập ký âm nốt huấn luyện trên nhạc cụ & Cho đoạn
nhạc cụ; có thể nhảy quãng tám trên giọng trần; bỏ qua bộ phân đoạn \\
\texttt{pyin} & $f_0$ xử lý tín hiệu cổ điển cộng bộ phân đoạn của đồ án & ``da-da-da''
staccato rõ nét; mặc định thư viện \\
\bottomrule
\end{tabular}
\end{table}

\subsection{Nhận biết lưới: phân đoạn và hợp nhất}

Ý tưởng phương pháp quan trọng nhất trong đường ống là phân đoạn \emph{nhận biết
lưới}. Bộ phân đoạn được truyền nhịp độ đã biết và một đường bao năng lượng cửa sổ
ngắn, mịn. Một chỗ khép phụ âm ngắn hoặc nhẹ giữa hai nốt cùng cao độ hiện ra thành
một chỗ trũng hẹp trong đường bao mịn đó; một cổng bề rộng tách một chỗ khép như
vậy khỏi một đợt tremolo rộng, và một chỗ trũng nông hơn vẫn được chấp nhận khi nó
rơi vào một phách, vì lưới cho biết một ranh giới nốt được kỳ vọng ở đó. Đây là thứ
tách các nốt cùng cao độ lặp lại, và nó thay thế một bộ tách thung lũng thô trước
đây vốn làm nhòe những chỗ khép ngắn.

Hợp nhất là lực cân bằng cần thiết. Vì vibrato và hơi thở có thể xé một nốt giữ hơi
thành nhiều mảnh, một bước hợp nhất chạy cho mọi backend nối các mảnh cùng cao độ
lại với nhau. Việc nhận biết lưới của nó là hàng rào bảo vệ: nó từ chối nối băng qua
một onset lưới, nên phân đoạn và hợp nhất không thể đấu nhau thành một dao động.

\subsection{Hiệu chỉnh quãng tám dựa trên phổ}
\label{sec:octave}

Một bước riêng, chèn giữa bộ theo dõi và cổng xác định tiếng cho các backend cấp
khung, sửa lỗi sóng hài phụ do Viterbi mô tả ở Mục~2. Bước này chạy cho mọi backend
và dựa trên một dấu hiệu phổ. Một tần số cơ bản thực tại $f$ có năng lượng ở cả các
hài lẻ ($f, 3f, 5f$) lẫn các hài chẵn ($2f, 4f, 6f$). Một lỗi sóng hài phụ, trong
đó bộ theo dõi báo cáo một nửa tần số thực, chỉ có năng lượng ở các bội chẵn, vì đó
mới là tần số cơ bản thực và các hài của nó; các bội lẻ của tần số được báo cáo có
ít hoặc không có năng lượng. Bước này vì thế tính, cho mỗi khung có tiếng, độ nổi
bật tại các bội lẻ và chẵn của $f_0$ được báo cáo, và nhân đôi tần số được báo cáo
chỉ khi độ nổi bật của hài lẻ sụp đổ so với độ nổi bật của hài chẵn. Phép thử được
thiết kế thận trọng: nó để yên các tần số cơ bản thực, kể cả các giọng thiếu tần số
cơ bản (missing-fundamental), vốn vẫn giữ năng lượng ở các hài lẻ.
Mục~\ref{sec:results-note} báo cáo tác động của bước này lên phần đánh giá.

\subsection{Trình soạn thảo tương tác (Chế độ Manual)}

Bản nháp tự động không phải lúc nào cũng đúng, và mắt xích yếu là phân đoạn chứ
không phải đường cao độ. Ứng dụng web vì thế cung cấp một chế độ Manual: một trình
soạn thảo dựa trên khuông nhạc, ngay trong trình duyệt, để sửa bản nháp tự động. Nó
gần như hoàn toàn ở phía máy khách. Nó hiện thực lại phần xây dựng bản nhạc và đánh
vần hòa âm của bước xuất ngay trong trình duyệt, khắc nhạc bằng một bản dựng
WebAssembly được nhúng của phần mềm ký hiệu nhạc Verovio~\cite{verovio}, và cung cấp
các thao tác chỉnh sửa thuần (đổi cao độ, đổi trường độ, gộp, tách, xóa thành lặng,
chèn) kèm hoàn tác và làm lại. Một dải cao độ tham chiếu hiển thị đường cong ngân
nga so với các nốt đã chọn, để người dùng thấy nơi bản nháp lệch khỏi thứ họ đã
hát. Máy khách chỉ gửi các nốt đã sửa về máy chủ ở các mép, để xuất lại MIDI và
MusicXML và tính lại giọng và hợp âm.

Chế độ Manual cũng chấp nhận các chỉnh sửa bằng ngôn ngữ tự nhiên. Một ô ``Ask''
nhận một chỉ dẫn như ``gộp các nốt 3 đến 5'' hay ``nâng hai nốt cuối lên một quãng
tám''; các cách diễn đạt phổ biến được khớp bởi một văn phạm ngoại tuyến chạy không
cần mạng, và mọi thứ khác được biến thành một kịch bản chỉnh sửa nhỏ bởi một mô hình
ngôn ngữ tùy chọn (Mục~\ref{sec:coaching}). Dù bằng cách nào, chỉ dẫn cũng trở thành
cùng một kịch bản đã được kiểm định, mà một trình thông dịch thuần trong trình duyệt
áp dụng như một bước hoàn tác được duy nhất, từ chối cả kịch bản nếu bất kỳ tham
chiếu nốt nào nằm ngoài phạm vi. Một nút ``Reorganise'' riêng đi hướng ngược lại:
thay vì một chỉnh sửa cục bộ có giới hạn, nó viết lại toàn bộ giai điệu thành các
giá trị nốt gọn, dễ đọc trong một bước hoàn tác được (Mục~\ref{sec:coaching}).

Các hợp âm dưới bản nháp có thể được làm lại từ thẻ kết quả, độc lập với việc chỉnh
nốt. Nhấp vào một hợp âm sẽ đưa ra các lựa chọn thay thế, tính ngoại tuyến: cho ô
nhịp đó, hệ thống xếp hạng các hợp âm ứng viên, trong một tập phong cách đã chọn gồm
hợp âm ba, hợp âm bảy nguyên bậc và át phụ (secondary dominant), theo mức mỗi hợp âm
phủ được bao nhiêu giai điệu của ô nhịp, kèm một khoản phạt độ phức tạp nhỏ để một
hợp âm ba đơn giản thắng khi hòa. Một nút ``Reharmonise'' thay vào đó hỏi mô hình
ngôn ngữ tùy chọn một tiến trình hợp âm hoàn toàn mới theo một phong cách yêu cầu
(``jazz hơn'', ``đơn giản hơn cho guitar''); mô hình đề xuất một hợp âm mỗi ô nhịp
từ một danh sách chất lượng cố định, và hệ thống kiểm định đề xuất, đánh vần và gán
nhãn số La Mã bằng music21~\cite{music21}, chấm độ khớp giai điệu từng ô nhịp để một
xung đột được hiển thị thay vì bị giấu, và khắc lại bản nhạc. Áp dụng kết quả nào
cũng cập nhật bản nhạc, phần phát lại, trình soạn thảo và bộ dịch giọng cùng lúc.

\subsection{Bộ công cụ mở rộng: phân tích, luyện tập, dịch giọng và luyện tai}
\label{sec:toolkit}

Ứng dụng web có nhiều thẻ. Ngoài Bộ ký âm mô tả ở trên, năm thẻ là các công cụ
riêng không đi qua bộ phân đoạn ký âm; cùng với bộ ký âm chúng tạo thành bộ công cụ
giới thiệu ở Mục~\ref{sec:toolkit-intro}. Mỗi cái tái sử dụng nền tảng chung, tức
phát hiện cao độ và ngăn xếp Web Audio cùng ký hiệu nhạc, cho một nhiệm vụ âm nhạc
khác. Bốn công cụ nhận đầu vào giọng dưới đây chia sẻ một quy tắc thiết kế với bộ ký
âm: logic thuần của chúng được tách khỏi DOM để kiểm thử đơn vị được, và micro được
giải phóng mỗi khi rời khỏi một thẻ.

\begin{itemize}
    \item \textbf{Bộ tìm cao độ (Pitch Finder)} phân tích một đoạn tải lên để tìm
    giọng, nhịp độ và một mã Camelot, dùng một phương pháp dựa trên chroma hoạt động
    được trên nhạc đa âm. Nó tái sử dụng mã chấm giọng và gấp nhịp độ nhưng không
    dùng bộ phân đoạn đơn âm, vốn sẽ sai với một bài hát đầy đủ.
    \item \textbf{Bộ thời gian thực (Realtime)} là một bộ giám sát cao độ trực tiếp
    và bộ lên dây đàn guitar ở phía máy khách xây trên Web Audio API, kèm một tập
    bài luyện giọng (một làn nốt mục tiêu, một trò chơi ghép nốt, một bộ luyện thang
    âm, phân tích vibrato, một bộ tìm quãng giọng và một nhật ký tiến bộ theo phiên).
    Một vòng đi về máy chủ không thể là thời gian thực, nên đường này chạy hoàn toàn
    trong trình duyệt và chỉ hỏi máy chủ giọng đã hát khi dừng, từ một biểu đồ tần
    suất cao độ, không tải lên âm thanh. Bộ tìm quãng giọng là bài duy nhất đưa ra
    một khẳng định thực nghiệm về người hát so với một quần thể: từ quãng phát âm cực
    đại đo được, nó báo cáo một phân vị (\emph{rộng hơn X\% người trưởng thành}). Để
    giữ khẳng định đó trung thực, trung bình của mô hình lấy từ một nghiên cứu được
    trích dẫn về quãng phát âm của người trưởng thành~\cite{hollien}, còn độ lệch
    chuẩn là một giả định được ghi rõ (khoảng bốn nửa cung, phù hợp với các tài liệu
    về hồ sơ quãng giọng) chứ không phải một giá trị đo được, và cả hai đều được in
    trong giao diện.
    \item \textbf{Bộ dịch giọng (Transposer)} chuyển một bản nhạc sang giọng mới.
    Chế độ tệp dịch một tệp MIDI hoặc MusicXML tải lên ở phía máy chủ bằng
    music21~\cite{music21}, di chuyển mọi bè và cả bộ khóa; chế độ ngân nga dịch bản
    ký âm gần nhất ở phía máy khách. Chế độ tệp cũng cung cấp các giọng tương thích
    Camelot dưới dạng các phím dịch một chạm.
    \item \textbf{Bộ hát theo (Sing-Along)} biến một bản nhạc MIDI hoặc MusicXML tải
    lên thành một bài tập karaoke: giai điệu được phát lại làm hướng dẫn trong khi
    giọng người dùng được theo dõi trực tiếp và chấm điểm từng nốt. Vì toàn bộ giai
    điệu đã biết trước, chính bản nhạc điều khiển đồng hồ; không có việc gióng thời
    gian tự do. Máy chủ rút gọn một tệp tải lên có thể đa âm thành một dòng giai điệu
    đơn âm duy nhất (đường viền cao nhất - skyline, bỏ các dấu nối), và trình duyệt
    lập lịch đếm vào và các nốt hướng dẫn tại các thời điểm tuyệt đối của đồng hồ âm
    thanh rồi chấm cao độ đã hát so với nốt đang hoạt động. Dung sai đúng cao độ có
    thể chọn (bốn mức độ khó), âm lượng hướng dẫn điều chỉnh được, và lượt hát có thể
    tạm dừng và tiếp tục.
    \item \textbf{Bộ luyện tai (Ear Trainer)} là tấm gương của các thẻ khác: thay vì
    người dùng tạo ra âm thanh để hệ thống phân tích, hệ thống phát ra âm thanh để
    người dùng nhận diện. Nó hoàn toàn ở phía máy khách, không micro và không máy
    chủ. Bốn chế độ con, chọn như các bài Realtime, phủ nhận diện quãng, chất lượng
    hợp âm, thang âm hoặc điệu thức, và bậc thang âm; mỗi câu hỏi phát chính các mẫu
    piano mà các thẻ khác dùng và đưa ra đáp án trắc nghiệm, kèm ba tập độ khó và một
    lịch sử độ chính xác theo từng chế độ giữ trong trình duyệt. Chế độ giọng kiểm
    tra nhận diện bậc thang âm sau một tiến trình kết I-IV-V-I, tức kỹ năng cao độ
    tương đối luyện được, chứ không phải gọi tên giọng tuyệt đối, vốn đòi hỏi thính
    giác tuyệt đối.
\end{itemize}

\subsection{Các tính năng AI tùy chọn và ngoại lệ với nguyên tắc cục bộ}
\label{sec:coaching}

Hệ thống cung cấp năm tính năng tùy chọn được hỗ trợ bởi một mô hình ngôn ngữ lớn
bên ngoài, và chúng là những nơi duy nhất mà bất kỳ dữ liệu nào rời khỏi máy. Chúng
chia sẻ một hợp đồng thiết kế, để ngoại lệ với nguyên tắc chỉ chạy cục bộ giữ được
nhỏ và đồng nhất. Mỗi tính năng chỉ kích hoạt khi có một cú nhấn nút rõ ràng; mỗi
tính năng chỉ gửi một bản tóm tắt văn bản hoặc số gọn, không bao giờ là âm thanh,
không bao giờ là bản ghi, thậm chí không cả tên tệp; và mỗi tính năng hiển thị phản
hồi của mô hình dưới dạng văn bản thuần chứ không phải mã đánh dấu, nên đầu ra không
tin cậy không thể chèn nội dung vào trang. Khóa API được đọc ở phía máy chủ, theo
từng yêu cầu, từ một tệp cấu hình cục bộ bị bỏ khỏi git và không bao giờ lộ cho
trình duyệt. Cả năm dùng chung một mô-đun truyền tải duy nhất, và tất cả đều tắt
theo mặc định và bất động cho tới khi một khóa được cấu hình, nên phần còn lại của
hệ thống vẫn hoàn toàn cục bộ.

Năm tính năng là:

\begin{itemize}
    \item \textbf{Gợi ý luyện tập cho Sing-Along.} Sau một lượt hát, máy khách trước
    tiên tính một phân tích phần trình bày, xác định và đọc được, ngay trong trình
    duyệt: độ lệch cao độ có dấu (nhất quán cao hay thấp), độ trôi từ đầu tới cuối
    lượt, các lần trượt quãng tám, độ chính xác của bước nhảy so với bước liền,
    quãng giọng yếu nhất, và các nốt tệ nhất kèm số ô nhịp. Phân tích này luôn được
    hiển thị và không cần mạng. Chỉ khi được yêu cầu thì bản tóm tắt số của nó, kèm
    ngữ cảnh âm nhạc, mới được gửi đi để lấy phản hồi bằng ngôn ngữ và gợi ý luyện
    tập.
    \item \textbf{Kế hoạch luyện tập.} Trên hub, báo cáo tiến bộ tổng hợp, tính cục
    bộ từ lịch sử luyện tập, có thể được gửi đi để lấy một kế hoạch một tuần với các
    mục nêu tên chính các bài của ứng dụng. Kế hoạch gần nhất được lưu đệm trong
    trình duyệt và gắn khóa theo một hàm băm của báo cáo, nên mở hub không bao giờ
    tự phát ra một yêu cầu.
    \item \textbf{Chỉnh sửa bằng ngôn ngữ tự nhiên.} Trong chế độ Manual, ô ``Ask''
    gửi một liệt kê văn bản của giai điệu và chỉ dẫn của người dùng, và mô hình trả
    về một kịch bản chỉnh sửa có cấu trúc nhỏ; trình duyệt kiểm định và áp dụng nó.
    Các chỉ dẫn phổ biến được xử lý hoàn toàn ngoại tuyến bằng một văn phạm và không
    bao giờ tới đường này.
    \item \textbf{Viết lại toàn bộ giai điệu.} Cũng trong chế độ Manual, một nút
    ``Reorganise'' gửi cao độ và độ dài theo phách của giai điệu hiện tại và nhận về
    một bản viết lại gọn hơn, với trường độ được gióng vào các giá trị nốt đơn giản
    và các mảnh vụn nối được gộp; trình duyệt kiểm định và áp dụng nó như một bước
    hoàn tác được duy nhất. Nơi ô Ask làm một chỉnh sửa cục bộ có giới hạn, tính năng
    này viết lại cả dòng cho dễ đọc.
    \item \textbf{Phối lại hợp âm (Reharmonisation).} Trên thẻ kết quả, một phong
    cách yêu cầu và hồ sơ lớp cao độ từng ô nhịp của giai điệu được gửi đi, và mô
    hình đề xuất một tiến trình hợp âm mới, mà hệ thống kiểm định, đánh vần, chấm
    điểm và khắc cục bộ.
\end{itemize}

Các yêu cầu phản hồi được xử lý ở phía máy chủ bởi các trình xử lý tuyến đồng bộ,
mà khung web chạy trên một luồng công nhân, nên một cuộc gọi tới một phút không chặn
vòng lặp sự kiện. Mô hình mặc định là một mô hình suy luận tiêu tốn các token ẩn
trước câu trả lời, nên các tính năng có cấu trúc (chỉnh sửa, viết lại và phối lại)
yêu cầu một ngân sách token lớn hơn các tính năng văn xuôi và có thể được trỏ tới
một mô hình nhanh hơn qua cấu hình. Ngoài năm cuộc gọi này, hoạt động của HumJob
hoàn toàn cục bộ.

\subsection{Phát hiện nhịp độ và điểm yếu đã biết của nó}
\label{sec:tempo-design}

Vì nhịp độ sai là nguyên nhân thường gặp của các mảnh nốt nối, ứng dụng cung cấp
một nút ``tìm nhịp độ của tôi''. Mục đích của nó hẹp: người dùng ngân nga vài phách
đều, hệ thống ước lượng một nhịp độ thoải mái, rồi người dùng thu lượt thật theo một
tiếng đếm ở nhịp độ đó, để lượng tử hóa rơi vào các phách nguyên.

Bộ ước lượng tính một đường bao độ mạnh onset bằng librosa~\cite{librosa}, rồi tìm
chu kỳ trội của đường bao đó bằng tự tương quan, có trọng số bởi một tiên nghiệm đặt
quanh một nhịp độ ngân nga thoải mái, và gấp kết quả vào một dải hợp lý theo quãng
tám. Tính chất quan trọng, cũng là một điểm yếu đã biết được xem xét ở
Mục~\ref{sec:bpm}, là nó trả về một \emph{nhịp độ toàn cục duy nhất} trung bình trên
cả đoạn, nghiêng về tiên nghiệm. Nó giả định một nhịp độ gần như không đổi và các
onset tương đối sạch, và nó dễ hỏng khi một trong hai giả định bị vi phạm.

Một hàng rào bảo vệ thứ hai, độc lập, nằm ngay trong bộ lượng tử hóa. Vì ngay cả một
lượt ngân nga cẩn thận cũng trôi vài phần trăm khỏi máy đếm nhịp, bộ lượng tử hóa
tinh chỉnh nhịp độ trước khi gióng: nó quét một dải có giới hạn quanh nhịp độ đã nêu
(khoảng $\pm 25\%$) và chọn độ dài phách làm cho nhịp điệu đã gióng đơn giản nhất về
mặt nhịp, để yên một lượt vốn đã đúng phách. Bản nhạc vẫn được ký âm ở nhịp độ người
dùng đặt; chỉ lưới mà onset gióng vào là theo nhịp độ thực hiện. Điều này hấp thụ
một lỗi nhịp độ vừa phải, dù từ người dùng hay từ bộ phát hiện trên, nhưng không hấp
thụ một lỗi lớn, đó là lý do sự dễ hỏng của bộ phát hiện vẫn quan trọng với các lỗi
lớn (Mục~\ref{sec:bpm}).

\subsection{Ngăn xếp công nghệ và cách chạy hệ thống}
\label{sec:stack}

Bảng~\ref{tab:stack} liệt kê các thành phần chính. Toàn bộ hệ thống chạy cục bộ.
Ứng dụng web là một máy chủ FastAPI phục vụ một giao diện trình duyệt tĩnh và tải
lại khi có chỉnh sửa mã nguồn; công cụ dòng lệnh ký âm một tệp trực tiếp.

\begin{table}[H]
\centering
\caption{Ngăn xếp công nghệ.}
\label{tab:stack}
\begin{tabular}{l p{0.60\textwidth}}
\toprule
\textbf{Lớp} & \textbf{Công nghệ} \\
\midrule
Đường ống lõi     & Python 3.12, NumPy \\
Backend cao độ    & pYIN qua librosa~\cite{librosa}; torchcrepe~\cite{crepe};
pesto-pitch~\cite{pesto}; penn / FCNF0++~\cite{penn}; basic-pitch trên ONNX
Runtime~\cite{basicpitch} \\
Onset / nhịp độ / chroma & librosa~\cite{librosa} \\
Mô hình bản nhạc, MIDI, MusicXML & music21~\cite{music21} \\
Khắc nhạc         & Verovio~\cite{verovio} (phía máy chủ, và một bản dựng
WebAssembly được nhúng trong trình duyệt) \\
Backend web       & FastAPI, Uvicorn (kèm tự động tải lại) \\
Giao diện web     & HTML / CSS / JavaScript tĩnh, Web Audio API \\
Tính năng AI (tùy chọn) & Máy khách httpx dùng chung tới một API mô hình ngôn ngữ
bên ngoài, tương thích OpenAI (DeepSeek); cấp năng lượng cho gợi ý luyện tập, kế
hoạch luyện tập, chỉnh sửa bằng ngôn ngữ tự nhiên, viết lại toàn bộ giai điệu và
phối lại hợp âm; tắt theo mặc định, tùy chọn bật, chỉ gửi tóm tắt văn bản hoặc số \\
Đánh giá          & mir\_eval~\cite{mireval} cho các chỉ số nốt \\
\bottomrule
\end{tabular}
\end{table}

% =========================================================================
% MUC 4: DANH GIA VA KET QUA
% =========================================================================
\section{Đánh giá và kết quả}

\subsection{Phương pháp đánh giá}

Mục này đánh giá bộ ký âm, lõi kỹ thuật của bộ công cụ và là thành phần duy nhất ước
lượng một nhãn chuẩn ẩn. Năm công cụ còn lại là các sản phẩm chức năng chứ không
phải đối tượng đo lường: logic thuần của chúng (việc xây bản nhạc và đánh vần hòa âm
ở máy khách, số học giọng của bộ dịch giọng, chấm điểm sing-along, và sinh câu hỏi
luyện tai) được kiểm thử đơn vị, và Phụ lục~A minh họa từng cái đang chạy. Không có
mốc chuẩn độ chính xác cho chúng vì chúng là tiện ích, không phải bộ ước lượng, nên
các phép đo dưới đây chỉ liên quan đến đường ống ký âm.

Phép đánh giá tiêu chuẩn của đồ án cho bộ ký âm là một khung dữ liệu tổng hợp có
nhãn. Nó dựng một tập giai điệu đã biết thành âm thanh, chạy đường ống một lần cho
mỗi mẫu, và so sánh đầu ra với tham chiếu đã biết. Các mẫu phủ tám giai điệu: một
thang âm Đô trưởng, một thang âm La thứ, một hợp âm rải, các nốt lặp, một giai điệu
có khoảng lặng, các bước nhảy quãng tám, một đoạn trích Twinkle, và một câu nhịp
điệu hỗn hợp.

Khung đo chấm hai thứ khác nhau, vì chúng đo các lỗi khác nhau:

\begin{itemize}
    \item \textbf{Phát hiện nốt (F1 của cao độ và onset)}, tính bằng
    mir\_eval~\cite{mireval}. Một nốt dự đoán khớp với một nốt tham chiếu khi cao độ
    đúng và onset nằm trong một dung sai; trường độ bị bỏ qua và không áp dụng dung
    sai bù trường độ. Việc này chạy ở hai dạng: một lượt \emph{sạch} mà cổng hồi quy
    kỳ vọng vượt qua với F1 $=1.0$, và một lượt \emph{thực tế} với vibrato rộng,
    tremolo, trôi cao độ, độ dài khoảng phụ âm dao động và các chỗ khép phụ âm không
    trọn.
    \item \textbf{Nhịp điệu}, thêm vào sau, vì F1 của nốt mù nhịp điệu. Nó chấm
    onset và trường độ đã lượng tử hóa của mỗi nốt so với một lưới dự kiến và được
    mô tả ở Mục~\ref{sec:results-rhythm}.
\end{itemize}

Cả hai đều là điểm tổng hợp, và cả hai đều bão hòa ở mức trần trong quá trình làm đồ
án (Mục~\ref{sec:results-note} và~\ref{sec:results-rhythm}). Vì một con số đã bão
hòa không thể dẫn dắt cải tiến và không nói gì về ngân nga thực, phần đánh giá sau
đó được mở rộng thành bốn loại bằng chứng, trình bày ở Mục~\ref{sec:four-evidence},
mỗi loại được thiết kế để không cần một người trình bày chính xác: các tỷ lệ lỗi
chẩn đoán trên một kho ngữ liệu đa dạng, các tính chất đúng đắn kiểu biến hình, một
bước tái tạo khứ hồi trên ngân nga thực, và một hiệu chuẩn độ thực tế của dữ liệu
tổng hợp so với các đoạn thu thực. Các kết quả F1 đã bão hòa dưới đây là mốc cơ sở
đã thúc đẩy việc mở rộng đó.

Một ghi chú về backend mà khung đo chạy. Các phép đo tổng hợp (F1 của nốt và nhịp
điệu, kho ngữ liệu chẩn đoán, và các tính chất biến hình tổng hợp) chạy backend
pyin xác định: nó không cần tải mô hình và tái lập được, và vì đóng góp nghiên cứu
là bộ phân đoạn và bộ lượng tử hóa chứ không phải mô hình cao độ (RQ2), chạy bộ theo
dõi đơn giản nhất cô lập được các bước đang xét. Các bằng chứng chạy trên bốn đoạn
ngân nga thực đã thu (tái tạo khứ hồi, hiệu chuẩn độ thực tế, và các phép kiểm biến
hình trên âm thanh thực) dùng PESTO thay thế, tức mặc định của ứng dụng web, để một
nửa phần đánh giá dùng đúng đường mà một người dùng thực nhận được. Các con số F1
tổng hợp dưới đây vì thế là kết quả của pyin và không nên đọc như độ chính xác của
PESTO trên giọng thực.

\subsection{Bài toán nhãn chuẩn}
\label{sec:groundtruth}

Trước các kết quả, cần nêu rõ giới hạn trung tâm của chúng. \textbf{Tất cả nhãn
chuẩn dùng trong báo cáo này đều là tổng hợp.} Các mẫu là những giai điệu mà chính
đồ án dựng thành âm thanh, nên cao độ, onset và trường độ dự kiến của chúng được biết
chính xác theo cách dựng. Đây là một điểm mạnh thực cho kiểm thử hồi quy, nhưng là
một điểm yếu thực cho tính hiệu lực: một đoạn ngân nga tổng hợp đều đặn hơn một đoạn
thực, và một mốc chuẩn chỉ xây trên ngân nga tổng hợp có thể bão hòa trong khi hệ
thống vẫn hỏng trên đầu vào thực.

Đồ án không có kho ngữ liệu ngân nga thực có nhãn. Lý do đặc thù cho nhiệm vụ này và
cũng là kiểu vấn đề giới hạn mọi mốc chuẩn ký âm: lấy nhãn chuẩn cho một đoạn ngân
nga thực đòi hỏi biết chính xác cao độ, onset và trường độ của mọi nốt mà người đó
dự định, và một người không thể tin cậy dán nhãn cao độ của chính đoạn ngân nga của
mình bằng tai. Kết quả là mọi con số phát hiện nốt và nhịp điệu trong báo cáo này đo
mức khớp với nhãn chuẩn \emph{tổng hợp} và phải đọc như một tín hiệu hồi quy, không
phải một thước đo độ chính xác thực tế. Mục~\ref{sec:eval-limits} đề xuất một biện
pháp cụ thể.

\subsection{Kết quả phát hiện nốt}
\label{sec:results-note}

Trên các mẫu sạch, đường ống đạt F1 nốt trung bình $1.000$. Con số nhiều thông tin
hơn là lượt thực tế, nơi bộ phân đoạn từng chật vật.
Bảng~\ref{tab:note-progression} ghi lại F1 thực tế trung bình qua ba vòng cải tiến.
Vòng đầu nâng nó từ $0.799$ lên $0.931$ bằng cách cho phân đoạn và hợp nhất nhận
biết lưới; vòng hai nâng lên $1.000$ bằng cách thêm hiệu chỉnh quãng tám dựa trên
phổ ở Mục~\ref{sec:octave}, sửa mẫu hỏng cuối cùng (các bước nhảy quãng tám, trước
đó trả về thấp một quãng tám).

\begin{table}[H]
\centering
\caption{F1 nốt trung bình trên các mẫu thực tế qua ba vòng. Lượt thực tế là lượt
giàu biểu cảm (vibrato rộng, tremolo, trôi, khoảng phụ âm dao động, chỗ khép không
trọn).}
\label{tab:note-progression}
\begin{tabular}{l l l}
\toprule
\textbf{Vòng} & \textbf{Thay đổi} & \textbf{F1 thực tế trung bình} \\
\midrule
Cơ sở                 & bộ tách thung lũng thô                 & $0.799$ \\
Nhận biết lưới        & phân đoạn + hợp nhất nhận biết lưới     & $0.931$ \\
Hiệu chỉnh quãng tám  & sửa sóng hài phụ dựa trên phổ           & $1.000$ \\
\bottomrule
\end{tabular}
\end{table}

Bảng~\ref{tab:note-perfixture} cho các chỉ số nốt từng mẫu trên lượt thực tế ở trạng
thái hiện tại. Mọi mẫu đạt $1.000$ về độ chính xác (precision), độ thu hồi (recall)
và F1, và giọng được phát hiện đúng trong mọi trường hợp. Đây là sự bão hòa được nói
đến xuyên suốt báo cáo này: nửa phát hiện nốt của khung tổng hợp không còn phân biệt
được một bộ phân đoạn tốt với một bộ kém hơn, nên tự nó không thể dẫn dắt cải tiến
nữa.

\begin{table}[H]
\centering
\caption{Chỉ số nốt từng mẫu trên lượt thực tế, trạng thái hiện tại. P, R, F1 là
độ chính xác, độ thu hồi và F1 từ mir\_eval; ref/est là số nốt tham chiếu so với số
nốt phát hiện.}
\label{tab:note-perfixture}
\begin{tabular}{l l l l l l}
\toprule
\textbf{Mẫu} & \textbf{P} & \textbf{R} & \textbf{F1} & \textbf{Giọng} &
\textbf{ref/est} \\
\midrule
c\_major\_scale  & $1.00$ & $1.00$ & $1.00$ & Đô trưởng & 8/8 \\
a\_minor\_scale  & $1.00$ & $1.00$ & $1.00$ & La thứ & 8/8 \\
arpeggio         & $1.00$ & $1.00$ & $1.00$ & Đô trưởng & 7/7 \\
repeated\_notes  & $1.00$ & $1.00$ & $1.00$ & Đô trưởng & 5/5 \\
with\_silence    & $1.00$ & $1.00$ & $1.00$ & Đô trưởng & 5/5 \\
octave\_leaps    & $1.00$ & $1.00$ & $1.00$ & Đô trưởng & 5/5 \\
twinkle          & $1.00$ & $1.00$ & $1.00$ & Đô trưởng & 14/14 \\
mixed\_rhythm    & $1.00$ & $1.00$ & $1.00$ & Đô trưởng & 5/5 \\
\midrule
\textbf{trung bình}    &        &        & $1.000$ &        & \\
\bottomrule
\end{tabular}
\end{table}

\subsection{Đánh giá nhịp điệu}
\label{sec:results-rhythm}

F1 nốt ở trên mù nhịp điệu theo bản chất: nó bỏ qua trường độ và chấp nhận bất kỳ
onset nào trong một dung sai, nên nó không nói gì về việc nhịp điệu ghi ra có đúng
không. Điều này quan trọng vì hai lỗi được báo cáo thường nhất trên ngân nga thực là
sai nhịp điệu (trường độ và onset ở sai phách) và chia nhỏ quá mức (một nốt trả về
thành các mảnh nối), và không lỗi nào nhìn thấy được với F1 nốt. Vì thế một chỉ số
nhịp điệu riêng được thêm vào. Nó so onset và trường độ đã lượng tử hóa của mỗi nốt
với một \emph{lưới dự kiến} đọc thẳng từ đặc tả của mỗi mẫu, và báo cáo tỷ lệ nốt có
onset đúng, trường độ đúng, và cả hai. Con số chủ đạo là ``cả hai'': đúng phách và
đúng độ dài ghi ra.

Hai điều kiện được chạy, chọn để khớp với hai thứ thực sự phá hỏng ngân nga thực.
Bảng~\ref{tab:rhythm} cho kết quả từng mẫu.

\begin{table}[H]
\centering
\caption{Độ chính xác nhịp điệu dưới hai điều kiện. ``on'', ``dur'' và ``both'' là
tỷ lệ nốt có onset đúng, trường độ đúng và cả hai; ``err'' là sai số onset trung
bình theo mili giây. Điều kiện dao động áp dao động thời gian của con người
($\pm 30$ ms) ở nhịp độ đúng; điều kiện sai nhịp độ áp một lỗi nhịp độ $+5\%$ lên
thời gian đúng lưới.}
\label{tab:rhythm}
\begin{tabular}{l r r r r @{\hskip 1.5em} r r r r}
\toprule
 & \multicolumn{4}{c}{\textbf{Dao động thời gian ($\pm 30$ ms)}}
 & \multicolumn{4}{c}{\textbf{Sai nhịp độ ($+5\%$)}} \\
\cmidrule(lr){2-5}\cmidrule(lr){6-9}
\textbf{Mẫu} & \textbf{on} & \textbf{dur} & \textbf{both} & \textbf{err} &
\textbf{on} & \textbf{dur} & \textbf{both} & \textbf{err} \\
\midrule
c\_major\_scale  & $1.00$ & $1.00$ & $1.00$ & $0$ & $1.00$ & $1.00$ & $1.00$ & $0$ \\
a\_minor\_scale  & $1.00$ & $1.00$ & $1.00$ & $0$ & $1.00$ & $1.00$ & $1.00$ & $0$ \\
arpeggio         & $1.00$ & $1.00$ & $1.00$ & $0$ & $1.00$ & $1.00$ & $1.00$ & $0$ \\
repeated\_notes  & $1.00$ & $1.00$ & $1.00$ & $0$ & $1.00$ & $1.00$ & $1.00$ & $0$ \\
with\_silence    & $1.00$ & $0.80$ & $0.80$ & $0$ & $1.00$ & $1.00$ & $1.00$ & $0$ \\
octave\_leaps    & $1.00$ & $1.00$ & $1.00$ & $0$ & $1.00$ & $1.00$ & $1.00$ & $0$ \\
twinkle          & $1.00$ & $1.00$ & $1.00$ & $0$ & $1.00$ & $1.00$ & $1.00$ & $0$ \\
mixed\_rhythm    & $1.00$ & $1.00$ & $1.00$ & $0$ & $1.00$ & $1.00$ & $1.00$ & $0$ \\
\midrule
\textbf{trung bình}    &        &        & $\mathbf{0.975}$ & & & & $\mathbf{1.000}$ & \\
\bottomrule
\end{tabular}
\end{table}

Hai phát hiện nổi bật. Thứ nhất, đúng lưới ở nhịp độ đúng thì bộ lượng tử hóa hoàn
hảo ($1.000$), và dưới dao động thời gian của con người trung bình giữ ở $0.975$:
mọi mẫu trừ một đạt điểm tuyệt đối, và chỉ mẫu có khoảng lặng mất một trường độ ghi
ra, nên bộ lượng tử hóa bền vững với độ trôi thời gian của một phần trình bày thực.
Thứ hai, một lỗi nhịp độ $5\%$ cũng được khôi phục về $1.000$, vì bộ lượng tử hóa
tinh chỉnh nhịp độ trước khi gióng (Mục~\ref{sec:tempo-design}): nó quét một dải có
giới hạn quanh nhịp độ đã nêu và chọn độ dài phách làm cho nhịp điệu đã gióng đơn
giản nhất về mặt nhịp, nên một đoạn ngân nga trôi vài phần trăm khỏi máy đếm nhịp
được kéo về các phách nguyên thay vì thành các onset lệch phách mà ký âm hiển thị
thành mảnh nối. Trong dải đó, nửa nhịp điệu của khung tổng hợp vì thế cũng bão hòa.

Hình~\ref{fig:rhythm} vẽ hai điều kiện cạnh nhau.

\begin{figure}[H]
\centering
\reportfig{0.85}{eval-rhythm.png}{Độ chính xác ``both'' của nhịp điệu từng mẫu dưới
dao động thời gian của con người (nhịp độ đúng) và dưới một lỗi nhịp độ $5\%$. Cả
hai điều kiện đều ở hoặc gần mức trần: bước tinh chỉnh nhịp độ hấp thụ lỗi nhịp độ,
và chỉ mẫu có khoảng lặng giảm nhẹ dưới dao động.}
\caption{Độ chính xác nhịp điệu dưới dao động thời gian và dưới nhịp độ sai.}
\label{fig:rhythm}
\end{figure}

\subsection{Vượt qua một con số đã bão hòa: bốn loại bằng chứng}
\label{sec:four-evidence}

Một chỉ số F1 ghim ở mức trần không thể tách một bộ phân đoạn tốt với một bộ kém
hơn, và nó hoàn toàn tổng hợp, nên nó không chứng minh gì về ngân nga thực. Cách sửa
thông thường, một kho ngữ liệu ngân nga thực có nhãn, bị chặn ở đây vì lý do ở
Mục~\ref{sec:groundtruth}: tác giả không phải là một ca sĩ chính xác và không tồn
tại tập dữ liệu ngân nga thực có nhãn. Vì thế phần đánh giá được mở rộng thành bốn
loại bằng chứng, mỗi loại chọn sao cho không cần một người trình bày chính xác, và
mỗi loại được báo cáo dưới đây kèm một tuyên bố rõ ràng về điều nó chứng minh và
không chứng minh. Cùng nhau chúng thay một con số đã bão hòa bằng một bức tranh vẫn
còn hữu ích.

\subsubsection{Bằng chứng 1: tỷ lệ lỗi chẩn đoán trên một kho ngữ liệu đa dạng}
\label{sec:diag}

Bước đầu là ngừng giấu các lỗi bên trong một con số tổng hợp. Thay vì F1 nốt, đường
ống được chấm trên một tập tỷ lệ lỗi chẩn đoán, mỗi tỷ lệ chỉ ra một dạng lỗi, trên
một kho ngữ liệu lớn hơn, đa dạng hơn. Kho có mười sáu giai điệu: tám mẫu gốc, ba
bản dân ca thuộc phạm vi công cộng (Ode to Joy, Mary Had a Little Lamb, Frere
Jacques), và năm giai điệu sinh theo thủ tục (bước ngẫu nhiên trong giọng, một hợp
âm rải, một bài tập quãng, và một dòng nhịp điệu biến đổi), tất cả đều dựng với nhãn
chuẩn đã biết. Mỗi giai điệu được dựng ở ba hồ sơ độ thực tế: \emph{sạch},
\emph{thực tế}, và một hồ sơ \emph{khó} mới, gay gắt hơn, thêm rung biên độ
(shimmer), rè cao độ (creak), một nền nhiễu tiếng thở, một tiếng vọng phòng ngắn, và
nhiễu micro hồng lên trên vibrato và các chỗ khép không trọn của hồ sơ thực tế.

Các tỷ lệ lỗi là: tỷ lệ chia thừa/thiếu có dấu (số nốt phát hiện trừ số nốt tham
chiếu, chia cho số tham chiếu, nên giá trị dương là chia nhỏ quá mức còn giá trị âm
là gộp hoặc bỏ nốt); tỷ lệ lỗi quãng tám (các nốt gióng được lệch đúng một quãng
tám); độ thu hồi lặp cùng cao độ (trong các cặp cùng cao độ liền kề của tham chiếu,
tỷ lệ còn sống sót thành hai nốt riêng, chính là than phiền ``gộp lặp''); và tỷ lệ
lệch lưới (các nốt có onset hoặc trường độ đã lượng tử hóa rời khỏi lưới nốt đơn
giản, chính là than phiền ``mảnh nối''). Bảng~\ref{tab:diag} cho các trung bình kho
theo hồ sơ.

\begin{table}[H]
\centering
\caption{Tỷ lệ lỗi chẩn đoán, trung bình kho trên mười sáu giai điệu, ở ba hồ sơ độ
thực tế. F1 là F1 nốt trung bình để tham chiếu. ``split'' là tỷ lệ chia thừa/thiếu
có dấu (âm là gộp hoặc bỏ). ``octave'' là tỷ lệ lỗi quãng tám. ``repeat'' là độ thu
hồi lặp cùng cao độ (cao hơn là tốt hơn). ``off-grid'' là tỷ lệ mảnh nối / lệch
lưới. Không giai điệu nào sụp về không nốt ở bất kỳ hồ sơ nào.}
\label{tab:diag}
\begin{tabular}{l r r r r r}
\toprule
\textbf{Hồ sơ} & \textbf{F1} & \textbf{split} & \textbf{octave} &
\textbf{repeat} & \textbf{off-grid} \\
\midrule
Sạch       & $1.000$ & $+0.00$ & $0.00$ & $1.00$ & $0.00$ \\
Thực tế    & $0.979$ & $-0.04$ & $0.00$ & $0.61$ & $0.00$ \\
Khó        & $0.918$ & $-0.04$ & $0.00$ & $0.63$ & $0.00$ \\
\bottomrule
\end{tabular}
\end{table}

Ba điều nổi bật. Thứ nhất, kho ngữ liệu phân biệt trở lại: các con số không còn toàn
$1.000$, và chúng xếp ba hồ sơ như kỳ vọng (sạch tốt nhất, khó tệ nhất về F1). Thứ
hai, con số đơn hữu ích nhất là độ thu hồi lặp, giảm từ $1.00$ trên đầu vào sạch còn
khoảng $0.61$ trên đầu vào thực tế: trên một lượt biểu cảm, bộ phân đoạn gộp khoảng
bốn trên mười cặp lặp cùng cao độ đáng lẽ phải giữ tách. Đó là điểm yếu ``gộp lặp''
từ lâu bị nghi ngờ, nay được đo chứ không chỉ báo cáo không chính thức. Thứ ba, tỷ
lệ quãng tám giữ ở không (hiệu chỉnh quãng tám dựa trên phổ giữ vững) và tỷ lệ lệch
lưới giữ ở không ở nhịp độ đúng (bộ lượng tử hóa giữ nốt ở các giá trị đơn giản).

Một kỳ vọng không được đáp ứng, và báo cáo nêu rõ điều đó. Hồ sơ \emph{khó} được
thiết kế để kích hoạt \emph{chia nhỏ quá mức} (một nốt giữ hơi bị vỡ thành mảnh
vụn), dạng lỗi mà bộ tổng hợp quá đều đặn trước đây không thể thể hiện. Phép đo trực
tiếp cho thấy điều ngược lại. Đường ống này không chia nhỏ quá mức ngay cả dưới nhiễu
loạn mạnh; các bước làm mượt bằng trung vị và hợp nhất giữ vững, và đầu vào khó thay
vào đó suy giảm bằng cách \emph{gộp hoặc bỏ} nốt. Đẩy bộ tổng hợp đủ xa để làm vỡ
một nốt thì lại làm nốt im bặt. Tỷ lệ chia có dấu vì thế âm nhẹ, không bao giờ dương.
Đây là một phát hiện thực về hành vi của đường ống, và nó sửa lại kỳ vọng trước đó
ghi trong báo cáo này.

\textbf{Điều này chứng minh:} hành vi của đường ống trên đầu vào đa dạng, biểu cảm
nay được đo qua các dạng lỗi riêng biệt, và kho ngữ liệu lại có thể phân biệt một bộ
phân đoạn tốt với một bộ kém hơn. Dạng lỗi chủ đạo với đầu vào biểu cảm được khu trú
về các nốt cùng cao độ bị gộp, không phải lỗi quãng tám hay nhịp điệu lệch lưới.
\textbf{Điều nó không chứng minh:} nhãn chuẩn vẫn là tổng hợp, nên các tỷ lệ tuyệt
đối phụ thuộc vào mức trung thực của các hồ sơ độ thực tế, chính là chủ đề của Bằng
chứng~4.

\subsubsection{Bằng chứng 2: các tính chất đúng đắn kiểu biến hình}
\label{sec:meta}

Loại bằng chứng thứ hai không cần giai điệu tham chiếu nào cả. Một phép kiểm biến
hình khẳng định một \emph{quan hệ} phải giữ giữa một lượt và một phiên bản đã biến
đổi của cùng lượt đó, nên ``nhãn'' là phép biến đổi, không phải một cao độ đã biết.
Năm tính chất được kiểm tra: bất biến dịch giọng (dịch cao độ đi $k$ nửa cung sẽ
dịch mọi nốt đầu ra đi $k$ và giữ nguyên số nốt); bất biến khuếch đại (nhân biên độ
để lại bản ký âm y hệt); bền vững với nhiễu cộng (thêm nhiễu trắng xuống tới tỷ số
tín hiệu trên nhiễu $30$~dB để lại các nốt không đổi); bất biến đệm khoảng lặng
(thêm khoảng lặng đầu hoặc cuối để lại cùng các nốt ở cùng giá trị); và bất biến
nhịp độ, hay bất biến co giãn thời gian (trình bày cùng giai điệu ở một nhịp độ đã
co giãn và đọc nó ở nhịp độ đó để lại các giá trị nốt không đổi). Cả năm đều giữ
trên các mẫu tổng hợp. Bất biến dịch giọng và bất biến khuếch đại cũng giữ trên các
đoạn thu thực, nơi chúng đúng bất kể đoạn ngân nga chính xác tới đâu.

\textbf{Điều này chứng minh:} một tập các quan hệ đúng đắn mà bộ ký âm phải tuân
theo, kiểm chứng được không cần một người trình bày chính xác và một phần trên cả âm
thanh thực; một thay đổi làm hỏng việc xử lý dịch giọng hay nhịp độ sẽ bị bắt ở đây.
\textbf{Điều nó không chứng minh:} một tính bất biến có thể giữ trong khi bản ký âm
vẫn sai theo một cách mà phép biến đổi cũng bảo toàn, nên chúng chặn được cả lớp lỗi
nhưng không đo độ chính xác tuyệt đối.

\subsubsection{Bằng chứng 3: một bước tái tạo khứ hồi trên ngân nga thực}
\label{sec:roundtrip}

Loại bằng chứng thứ ba là loại duy nhất tính trên giọng thực. Vì một đoạn ngân nga
thực không có nhãn chuẩn, bản ký âm được chấm so với \emph{chính âm thanh của nó}:
các nốt đã ký âm được tổng hợp lại thành một dải âm đơn giản, và bản tái tạo đó được
so với bản thu gốc. Mức khớp cao độ là tỷ lệ khung có tiếng ở cả hai nằm trong $50$
cent; mức khớp onset là một F1 thời điểm onset giữa các onset phát hiện của âm thanh
và các onset đã ký âm; và một tóm tắt ``giải thích được âm thanh'' duy nhất lấy
trung bình hai thứ đó. Bảng~\ref{tab:roundtrip} cho kết quả trên bốn đoạn thu thực.

\begin{table}[H]
\centering
\caption{Tái tạo khứ hồi trên bốn đoạn thu thực. ``explains'' là điểm tóm tắt,
``pitch'' là tỷ lệ khung có tiếng ở cả hai nằm trong $50$ cent, ``onset'' là F1
thời điểm onset, và ``cents'' là sai số cao độ tuyệt đối trung vị trên các khung có
tiếng ở cả hai. Cao hơn là tốt hơn trừ cents.}
\label{tab:roundtrip}
\begin{tabular}{l r r r r}
\toprule
\textbf{Đoạn thu} & \textbf{explains} & \textbf{pitch} & \textbf{onset} &
\textbf{cents} \\
\midrule
mary\_1     & $0.59$ & $0.70$ & $0.49$ & $40$ \\
repeated\_1 & $0.75$ & $0.92$ & $0.57$ & $20$ \\
scale\_1    & $0.54$ & $0.73$ & $0.35$ & $20$ \\
twinkle\_1  & $0.61$ & $0.84$ & $0.39$ & $20$ \\
\midrule
\textbf{trung bình} & $\mathbf{0.62}$ & & & \\
\bottomrule
\end{tabular}
\end{table}

Mức khớp cao độ cao (giữa $0.70$ và $0.92$), nghĩa là giai điệu tổng hợp lại bám
theo cao độ mà người đó thực sự đã ngân nga và không xảy ra trượt quãng tám lớn nào.
Mức khớp onset thấp hơn (giữa $0.35$ và $0.57$), phần lớn vì bộ phát hiện onset kích
hoạt ở nhiều chỗ trong đoạn ngân nga thô hơn số nốt thực có, nên con số hạ thấp mức
khớp thật. Quan trọng là, vì điểm số so mỗi bản ký âm với \emph{chính} âm thanh của
nó và không bao giờ với một giai điệu dự định, một đoạn ngân nga thiếu chính xác
không thể bị tính là lỗi của hệ thống: một nốt sai mà đường ống bắt được trung thành
thì vẫn khớp. Đó là điều cho phép những con số này đến từ một người ngân nga không
phải ca sĩ được huấn luyện. Số đoạn thu ít (bốn), nên chúng được đọc như một tín
hiệu kiểm tra, không phải một xu hướng.

Một đoạn thu đáng ghi chú riêng. Trên đoạn Mary-Had-a-Little-Lamb, đường ống đã
\emph{chia nhỏ quá mức} giai điệu (nó trả về khoảng gấp đôi số nốt dự định). Đó là
dạng lỗi mà hồ sơ tổng hợp \emph{khó} không thể tái tạo (Mục~\ref{sec:diag}), nên
giọng thực rõ ràng mang những gợi ý, hoặc nhiễu thu âm, mà bộ tổng hợp không có.
Đoạn đó cũng nhỏ tiếng và nhịp độ không chắc, nên nó là một điểm dữ liệu yếu duy
nhất chứ không phải một tỷ lệ đo được; nhưng nó là một lý do cụ thể để không đọc quá
mức phát hiện chỉ-tổng-hợp rằng đường ống này không chia nhỏ quá mức, và là thêm một
lập luận cho nhãn chuẩn thực.

\textbf{Điều này chứng minh:} một tín hiệu kiểm tra không giám sát, trên giọng thực,
không cần nhãn và bắt được sai lệch lớn về cao độ, onset và phân đoạn, kể cả trượt
quãng tám. \textbf{Điều nó không chứng minh:} nó là một tín hiệu kiểm tra, không
phải một con số về độ chính xác. Một lỗi tự nhất quán, như một độ lệch chỉnh dây
không đổi trong $50$ cent, hay bất kỳ lỗi nào mà bản tổng hợp lại tái tạo, sẽ để hai
đường cong khớp nhau và không bị bắt.

\subsubsection{Bằng chứng 4: hiệu chuẩn độ thực tế của dữ liệu tổng hợp từ các đoạn
thu thực}
\label{sec:calib}

Loại bằng chứng thứ tư khép lại vòng trở về loại thứ nhất. Các tỷ lệ lỗi chẩn đoán
của Bằng chứng~1 chỉ trung thực đúng bằng mức các hồ sơ độ thực tế mà chúng được đo
theo, nên các hồ sơ đó được kiểm tra so với giọng thực. Với mỗi đoạn thu thực, bộ
theo dõi được chạy và các đại lượng mà các núm độ thực tế mô hình hóa được đo: độ
sâu và tần số vibrato, độ trôi cao độ, dao động thời gian so với máy đếm nhịp, độ sâu
các chỗ trũng năng lượng phụ âm, và rung biên độ. Bảng~\ref{tab:calib} cho các trung
vị trên bốn đoạn thu thực, cạnh giá trị mà hồ sơ \emph{thực tế} giả định.

\begin{table}[H]
\centering
\caption{Các đại lượng độ thực tế đo từ bốn đoạn thu thực (trung vị), cạnh giá trị
mà hồ sơ tổng hợp thực tế giả định. Hàng dao động thời gian là ngoại lệ: hồ sơ cao
độ thực tế giữ thời gian sạch, nên $30$~ms hiển thị là dao động mà phép kiểm nhịp
điệu bơm vào (Mục~\ref{sec:results-rhythm}), chính là giá trị mà số đo được đem so.}
\label{tab:calib}
\begin{tabular}{l r r}
\toprule
\textbf{Đại lượng} & \textbf{Đo được (trung vị)} & \textbf{Hồ sơ thực tế} \\
\midrule
Độ sâu vibrato (cent) & $26$  & $55$ \\
Tần số vibrato (Hz)   & $5.2$ & $5.3$ \\
Độ trôi cao độ (cent) & $18$  & $18$ \\
Dao động thời gian (ms) & $131$ & $30$ \\
Độ sâu chỗ khép (dB)  & $5$   & $16$ \\
Rung biên độ (dB)     & $1.1$ & --- \\
\bottomrule
\end{tabular}
\end{table}

Tần số vibrato và độ trôi cao độ mà bộ tổng hợp giả định gần với những gì ngân nga
thực thể hiện: tần số khớp gần như chính xác ($5.2$ so với $5.3$~Hz), và độ trôi rơi
đúng giá trị giả định ($18$ cent ở cả hai). Độ sâu vibrato giả định rộng hơn số đo
được ($55$ so với $26$ cent), nên hồ sơ thực tế, nếu có gì, còn khó hơn một đoạn
ngân nga thực về cao độ. Dao động thời gian đo được lớn hơn nhiều giá trị dùng trong
phép kiểm nhịp điệu ($131$ so với $30$~ms), nghĩa là ngân nga thực lỏng lẻo về thời
gian hơn nhiều so với giả định của phép kiểm nhịp điệu và là một lý do khiến công
việc về nhịp độ ở Mục~\ref{sec:bpm} quan trọng.

\textbf{Điều này chứng minh:} độ thực tế của dữ liệu tổng hợp được đặt nền trên các
con số đo được thay vì đoán, và hai núm trung tâm (tần số vibrato và độ trôi) khớp
tốt với giọng thực. \textbf{Điều nó không chứng minh:} bốn đoạn thu là một mẫu nhỏ,
và các số đo là ước lượng từ bộ theo dõi, không phải chính xác.

\subsection{Diễn giải các điểm số}
\label{sec:interpret}

Các bằng chứng phải được đọc cùng nhau và đọc cẩn thận. Một F1 nốt bão hòa $1.000$
không có nghĩa hệ thống ký âm ngân nga thực đúng; nó có nghĩa mốc chuẩn phát hiện
nốt tổng hợp không còn phân biệt được một bộ phân đoạn tốt với một bộ tồi. Bốn loại
bằng chứng ở trên được xây dựng chính vì điều đó. Các tỷ lệ lỗi chẩn đoán (Bằng
chứng~1) phục hồi một mốc chuẩn biết phân biệt và khu trú dạng lỗi chủ đạo với đầu
vào biểu cảm về các nốt cùng cao độ bị gộp; các tính chất biến hình (Bằng chứng~2)
cho các quan hệ đúng đắn không cần người trình bày; bước tái tạo khứ hồi (Bằng
chứng~3) cho một tín hiệu giọng thực đầu tiên, không giám sát; và hiệu chuẩn độ thực
tế (Bằng chứng~4) giữ các hồ sơ tổng hợp trung thực. Không loại nào là một con số về
độ chính xác thực tế, vì điều đó vẫn cần nhãn chuẩn thực, nhưng cùng nhau chúng là
một bức tranh hữu ích hơn nhiều so với một con số ở mức trần.

Phần đánh giá nhịp điệu từng khu trú dạng lỗi thực tế chủ đạo về một nguyên nhân duy
nhất: sự nhạy cảm với lỗi nhịp độ. Chẩn đoán đó thúc đẩy bước tinh chỉnh nhịp độ,
nay hấp thụ một lỗi nhịp độ vừa phải và, cùng với nó, hầu hết hành vi ``sai nhịp
điệu'' mà một lệch nhịp độ nhỏ từng gây ra. Điều còn để ngỏ là hành vi của bộ phát
hiện nhịp độ trên ngân nga thực trôi vượt quá dải tinh chỉnh (Mục~\ref{sec:bpm}), và,
căn bản hơn, sự thiếu vắng bất kỳ nhãn chuẩn thực nào để đo độ chính xác tuyệt đối.

\subsection{Điểm yếu của phát hiện nhịp độ}
\label{sec:bpm}

Kết quả nhịp điệu ở trên cho thấy một lỗi nhịp độ vừa phải nay được hấp thụ bởi bước
tinh chỉnh trong bộ lượng tử hóa. Tuy nhiên, bước đó chỉ trải một dải có giới hạn
(khoảng $\pm 25\%$ nhịp độ đã nêu), nên một lỗi lớn từ chính bộ phát hiện nhịp độ
của hệ thống vẫn thoát khỏi nó, và đáng nêu tại sao bộ phát hiện đó dễ hỏng.

Như mô tả ở Mục~\ref{sec:tempo-design}, bộ phát hiện trả về một nhịp độ toàn cục duy
nhất: nó lấy tự tương quan của đường bao độ mạnh onset trên cả đoạn, đánh trọng số
bằng một tiên nghiệm đặt gần một nhịp độ ngân nga thoải mái, và trả về chu kỳ trội.
Điều này có ba hệ quả trên đầu vào thực. Thứ nhất, nó giả định một nhịp độ không
đổi; một đoạn ngân nga nhanh dần hay chậm dần có năng lượng onset ở nhiều hơn một
chu kỳ, và tự tương quan trả về một thỏa hiệp mờ thay vì một trong hai nhịp độ thực.
Thứ hai, tiên nghiệm chủ động thiên lệch ước lượng về tâm dải của nó, nên một đoạn
ngân nga thực sự chậm bị kéo lên. Thứ ba, nó dựa vào các onset sạch, cách đều, mà
ngân nga legato hoặc nhiễu ngoài đời không phải lúc nào cũng cung cấp. Trong thử
nghiệm không chính thức, một đoạn ngân nga giai điệu khoảng $76$ nhịp mỗi phút có bao
gồm một đoạn nhanh hơn được báo cáo là khoảng $101$, một giá trị đáng ngờ gần với
tiên nghiệm, chính là dấu hiệu tiên nghiệm áp đảo một tự tương quan yếu. Một lỗi cỡ
đó (khoảng $+33\%$) nằm ngoài dải tinh chỉnh nên không được khôi phục, đúng là
trường hợp mà sự dễ hỏng của bộ phát hiện vẫn cắn. Điểm yếu này là mục đầu của công
việc tương lai ở Mục~\ref{sec:directions}.

% =========================================================================
% MUC 5: HAN CHE VA CONG VIEC TUONG LAI
% =========================================================================
\section{Hạn chế và công việc tương lai}

\subsection{Hạn chế của đánh giá}
\label{sec:eval-limits}

Hạn chế chính là hạn chế nêu ở Mục~\ref{sec:groundtruth}: không tồn tại kho ngữ liệu
ngân nga thực có nhãn, nên không con số nào trong báo cáo này là một con số về độ
chính xác thực tế. Cả F1 nốt tổng hợp lẫn F1 nhịp điệu đều bão hòa trong quá trình
làm đồ án. Cách đáp là mở rộng phần đánh giá (Mục~\ref{sec:four-evidence}) thay vì
tiếp tục báo cáo một con số ở mức trần: các tỷ lệ lỗi chẩn đoán trên kho ngữ liệu đa
dạng phân biệt trở lại, các tính chất biến hình và bước tái tạo khứ hồi thêm các tín
hiệu không cần một người trình bày chính xác, và hiệu chuẩn độ thực tế giữ các hồ sơ
tổng hợp trung thực. Điều mà không loại nào cho là độ chính xác tuyệt đối trên giọng
thực, vốn vẫn cần nhãn chuẩn thực. Một phát hiện thứ hai thu hẹp kế hoạch trước đó:
lỗi chia nhỏ quá mức được kỳ vọng xuất hiện khi shimmer, tiếng thở và creak được
thêm vào bộ tổng hợp, nhưng phép đo trực tiếp cho thấy đường ống này không chia nhỏ
quá mức; các hàng rào làm mượt và hợp nhất của nó giữ vững, và đầu vào khó suy giảm
bằng cách gộp hoặc bỏ nốt. Vậy khoảng trống đánh giá còn mở không phải là tổng hợp
chia-nhỏ-quá-mức mà là nhãn chuẩn thực.

Biện pháp được đề xuất không phải là một tập tổng hợp lớn hơn mà là một tập nhỏ các
đoạn ngân nga thực, có nhãn. Trở ngại là nhãn chuẩn, vì một người không thể dán nhãn
cao độ của chính đoạn ngân nga của mình bằng tai. Kế hoạch, được tinh chỉnh trong
quá trình làm đồ án, đảo ngược việc dán nhãn. Thay vì ngân nga trước rồi dán nhãn
sau, nhãn được cố định trước khi thu: người dùng ngân nga một giai điệu họ biết,
theo một tiếng đếm ở nhịp độ đã biết, và nhãn chuẩn nhịp điệu đến từ bản nhạc đã
biết và lưới của tiếng đếm. Cao độ của mỗi nốt được xác nhận bằng một ứng dụng tìm
cao độ bên ngoài dùng như một công cụ hỗ trợ, với người dùng kiểm chứng từng nốt so
với bản nhạc họ chủ động ngân nga, thay vì tin công cụ ngoài như một lời sấm. Điều
này quan trọng vì một bộ tìm cao độ bên ngoài cũng chỉ là một bộ phát hiện cao độ
khác; nếu tin nó một cách mù quáng và nó tình cờ chia sẻ một lỗi với đường ống, phần
đánh giá sẽ chấm lỗi chung đó là đúng. Giữ người dùng trong vòng lặp, và lấy nhãn
nhịp điệu từ tiếng đếm chứ không phải từ công cụ, tránh được sự vòng vo đó. Các bản
thu đã kiểm chứng, với nốt và nhịp độ đã biết, sẽ cho đồ án nhãn chuẩn thực đầu
tiên.

\subsection{Các thiếu sót của hệ thống}

Một số thiếu sót của hệ thống hiện tại đã biết và được theo dõi.
Bảng~\ref{tab:deficiencies} liệt kê chúng cùng biện pháp dự kiến cho từng cái.

\begin{longtable}{>{\raggedright\arraybackslash}p{0.30\textwidth} >{\raggedright\arraybackslash}p{0.58\textwidth}}
\caption{Các thiếu sót đã biết của hệ thống hiện tại và công việc dự kiến.}
\label{tab:deficiencies} \\
\toprule
\textbf{Lĩnh vực} & \textbf{Mô tả và biện pháp dự kiến} \\
\midrule
\endfirsthead
\toprule
\textbf{Lĩnh vực} & \textbf{Mô tả và biện pháp dự kiến} \\
\midrule
\endhead
\midrule
\multicolumn{2}{r}{\footnotesize\itshape Tiếp tục ở trang sau.} \\
\endfoot
\bottomrule
\endlastfoot
Sự dễ hỏng của phát hiện nhịp độ & Bộ ước lượng ``tìm nhịp độ của tôi'' trả về một
nhịp độ toàn cục duy nhất bị thiên lệch bởi một tiên nghiệm, và không đáng tin trên
ngân nga thực có nhịp độ biến đổi, chậm, hoặc legato (Mục~\ref{sec:bpm}); bước tinh
chỉnh nhịp độ của bộ lượng tử hóa hấp thụ một lỗi vừa phải nhưng không hấp thụ một
lỗi lớn. Biện pháp dự kiến: làm phẳng hoặc bỏ tiên nghiệm, và dựng đường bao onset
từ chính các onset xác định tiếng của đường ống thay vì từ dòng phổ thô. \\
\addlinespace
Gộp các nốt cùng cao độ lặp lại & Trên đầu vào biểu cảm, bộ phân đoạn gộp khoảng bốn
trên mười cặp lặp cùng cao độ đáng lẽ phải giữ tách (độ thu hồi lặp khoảng $0.61$,
Mục~\ref{sec:diag}). Đây, chứ không phải chia nhỏ quá mức, là dạng lỗi chủ đạo với
đầu vào biểu cảm. Biện pháp dự kiến: tăng cường gợi ý onset cùng cao độ và đo lại độ
thu hồi lặp làm mục tiêu. \\
\addlinespace
Chia nhỏ quá mức không xảy ra, nhưng từng được kỳ vọng & Lỗi chia nhỏ quá mức (một
nốt trả về thành mảnh vụn) được kỳ vọng xuất hiện khi shimmer, tiếng thở và creak
được thêm vào bộ tổng hợp. Phép đo cho thấy đường ống này không chia nhỏ quá mức; nó
gộp hoặc bỏ nốt (Mục~\ref{sec:diag}). Không cần biện pháp; kỳ vọng được sửa lại. \\
\addlinespace
Không có nhãn chuẩn thực & Không con số nào là một con số về độ chính xác thực tế
(Mục~\ref{sec:groundtruth}); bốn đoạn thu thực chỉ nuôi bước tái tạo khứ hồi và hiệu
chuẩn, vốn khoan dung với người trình bày, không phải một điểm số độ chính xác. Biện
pháp dự kiến: quy trình thu ngân-nga-theo-bản-nhạc-đã-biết ở
Mục~\ref{sec:eval-limits}. \\
\addlinespace
Chỉ đơn âm & Đường ống ký âm giả định một giọng duy nhất, một nốt tại một thời điểm.
Đầu vào đa âm hoặc hợp âm nằm ngoài phạm vi; thẻ Pitch Finder dùng một phương pháp
dựa trên chroma riêng cho các bài hát đầy đủ. \\
\addlinespace
Tính sẵn có của backend tùy chọn & Các backend chính xác nhất (PESTO, FCNF0++) là
các cài đặt tùy chọn với cây phụ thuộc nặng hơn, và FCNF0++ tải trọng số khi dùng
lần đầu. Mặc định thư viện là pyin, không cần cài thêm; chọn một backend chưa cài
gây lỗi nhập (import) thay vì suy giảm âm thầm, nên một lần triển khai phải cài
backend nào được cấu hình (bản demo công khai dùng PESTO). \\
\addlinespace
Giả định nhịp cố định & Lượng tử hóa giả định một nhịp đơn giản và một nhịp độ đã
biết, ổn định. Thay đổi bộ chỉ nhịp và rubato biểu cảm không được mô hình hóa. \\
\end{longtable}

\subsection{Các hướng tương lai}
\label{sec:directions}

Ba hướng được ghi lại cho công việc sau báo cáo này, theo thứ tự ưu tiên.

\begin{itemize}
    \item \textbf{(A) Độ bền vững của nhịp độ (ứng viên hàng đầu).} Với bước tinh
    chỉnh nhịp độ của bộ lượng tử hóa nay đã hấp thụ các lỗi vừa phải, bộ phát hiện
    nhịp độ là đòn bẩy lớn nhất còn lại lên độ chính xác nhịp điệu thực tế cho các
    lỗi lớn nằm ngoài dải tinh chỉnh (Mục~\ref{sec:bpm}), và nó kiểm thử được trực
    tiếp vì các mẫu có nhịp độ đã biết. Hai thay đổi rẻ, tác động cao: làm phẳng
    hoặc bỏ tiên nghiệm nhịp độ để các nhịp độ chậm không bị kéo lên, và dựng đường
    bao onset từ chính các onset phụ âm của đường ống thay vì từ dòng phổ thô.
    \item \textbf{(B) Nhãn chuẩn ngân nga thực.} Cởi trói tối thượng, chỉ bị chặn ở
    phương pháp dán nhãn mà Mục~\ref{sec:eval-limits} giải quyết. Bốn đoạn thu thực
    đã nuôi bước tái tạo khứ hồi và hiệu chuẩn độ thực tế
    (Mục~\ref{sec:roundtrip} và~\ref{sec:calib}); một tập có nhãn lớn hơn, thu bằng
    cùng phương pháp, sẽ biến chúng từ một tín hiệu kiểm tra thành một tín hiệu chắc
    hơn và cho đồ án đo được độ chính xác thực tế lần đầu tiên.
    \item \textbf{(C) Gợi ý cho lặp bị gộp.} Với việc chia nhỏ quá mức đã được chứng
    minh không xảy ra, điểm yếu đo được là điểm ngược lại: các nốt cùng cao độ lặp
    lại bị gộp trên đầu vào biểu cảm (Mục~\ref{sec:diag}). Tăng cường gợi ý onset
    cùng cao độ, với độ thu hồi lặp làm chỉ số mục tiêu, là bước tiếp theo cụ thể về
    phân đoạn.
\end{itemize}

\subsection{Hướng tới mức sẵn sàng cho sản phẩm}

HumJob là một đồ án nghiên cứu và học thuật, không phải một dịch vụ sản phẩm. Nó
chạy cục bộ, một lựa chọn thiết kế có chủ đích, nhưng một lần triển khai rộng hơn sẽ
cần thêm: đóng gói các backend nơ-ron tùy chọn để chúng cài đặt đáng tin, một bộ
kiểm thử tự động rộng hơn trên các bản thu thực, và xử lý lỗi cho vô số cách mà một
bản thu micro thực có thể không dùng được (cắt đỉnh, nhiễu nền, quá nhỏ tiếng).
Không cái nào khó về mặt khái niệm, nhưng cùng nhau chúng là một khối lượng công việc
đáng kể vượt ngoài phạm vi của đồ án môn học này.

% =========================================================================
% MUC 6: KET LUAN
% =========================================================================
\section{Kết luận}

Đồ án này đặt mục tiêu ký âm một giai điệu ngân nga thành nốt, giọng và hợp âm, bằng
cách chấp nhận hai ràng buộc nhỏ về cách người dùng ngân nga và xây dựng toàn bộ
đường ống quanh chúng. Bộ ký âm thu được là một đường ống ưu tiên cục bộ gồm các hàm
thuần trên ba kiểu dữ liệu, với năm backend cao độ thay thế được cho nhau, một bộ
phân đoạn và hợp nhất nhận biết lưới, và một bước hiệu chỉnh quãng tám dựa trên phổ.
Quanh lõi đó, đồ án lớn dần thành một bộ công cụ ưu tiên cục bộ cho giọng hát, cung
cấp dưới dạng một ứng dụng web bảy thẻ: bộ ký âm và trình soạn thảo tương tác của
nó, một bộ tìm cao độ, một bộ dịch giọng, một bộ giám sát thời gian thực và luyện
giọng, một bộ hát theo, một bộ luyện tai, và một bảng luyện tập. Tất cả đều chạy
trên chính máy người dùng; ngoại lệ có chủ đích duy nhất là một họ nhỏ gồm năm tính
năng AI tùy chọn, chỉ văn bản (gợi ý luyện tập, một kế hoạch luyện tập, chỉnh sửa
bằng ngôn ngữ tự nhiên, viết lại toàn bộ giai điệu, và phối lại hợp âm), mỗi cái gửi
một bản tóm tắt gọn và không bao giờ là âm thanh.

Đối với các câu hỏi nghiên cứu, các kết luận trung thực là pha trộn. \textbf{RQ1}
được ủng hộ theo cách dựng và theo phần đánh giá: hai ràng buộc giao diện đúng là
biến các phần khó thành các phần khả thi, và các bước nhận biết lưới khai thác chúng
là thứ đã đưa F1 nốt thực tế từ $0.799$ lên $1.000$. \textbf{RQ2} được ủng hộ: các
lỗi chẩn đoán được là lỗi phân đoạn và lượng tử hóa trên một đường cao độ về cơ bản
đúng, không phải lỗi mô hình cao độ. Trường hợp rõ nhất là lỗi quãng tám, được chứng
minh là một hiện tượng do giải mã Viterbi chứ không phải một lỗi trong âm thanh, và
được sửa bằng một phép thử phổ chứ không phải bằng đổi mô hình. \textbf{RQ3} tạo ra
bài học quan trọng nhất của đồ án. Cả F1 nốt lẫn F1 nhịp điệu mà hệ thống được tinh
chỉnh theo đều bão hòa ở $1.000$ trong khi hệ thống vẫn chưa đáng tin trên ngân nga
thực, vì một đoạn ngân nga tổng hợp đều đặn hơn một đoạn thực và không có kho ngữ
liệu thực có nhãn để đo. Thay vì tiếp tục báo cáo một con số ở mức trần, phần đánh
giá được xây dựng lại quanh bốn loại bằng chứng không cần một ca sĩ chính xác
(Mục~\ref{sec:four-evidence}). Chúng phân biệt trở lại và, giữa chúng, nói lên điều
gì đó cụ thể: dạng lỗi chủ đạo với đầu vào biểu cảm là các nốt cùng cao độ bị gộp
(không phải, như kỳ vọng, chia nhỏ quá mức, điều mà đường ống này không làm); bộ ký
âm tuân theo một tập các bất biến đúng đắn cũng giữ trên âm thanh thực; và bản ký âm
của một đoạn ngân nga thực có thể được kiểm tra so với chính âm thanh của nó như một
tín hiệu kiểm tra không giám sát. Điều mà không loại nào cho là độ chính xác thực tế
tuyệt đối, vốn vẫn cần nhãn chuẩn thực.

Sản phẩm chính vì thế không phải là một con số độ chính xác nổi bật. Đó là một bộ
công cụ hoạt động được, có cấu trúc tốt, mà lõi ký âm được ghép với một phần đánh giá
trung thực nói rõ lõi đó đang ở đâu: mạnh trên mốc chuẩn tổng hợp tổng hợp, được
phân biệt trở lại bởi các chẩn đoán phong phú hơn, và chưa được kiểm chứng trên đầu
vào thực. Bốn đoạn thu thực đã neo các tín hiệu tái tạo khứ hồi và hiệu chuẩn độ
thực tế; bước giá trị nhất tiếp theo là mở rộng chúng thành một tập lớn hơn, đáng
tin, gồm các đoạn ngân nga thực có nhãn, để độ chính xác thực tế của hệ thống có thể
được đo trực tiếp thay vì suy ra từ một đại diện tổng hợp.

% =========================================================================
% TAI LIEU THAM KHAO
% =========================================================================
\newpage
\addcontentsline{toc}{section}{Tài liệu tham khảo}
\begin{thebibliography}{99}

\bibitem{pyin} M. Mauch and S. Dixon. \emph{pYIN: A fundamental frequency
estimator using probabilistic threshold distributions.} IEEE International
Conference on Acoustics, Speech and Signal Processing (ICASSP), 2014.

\bibitem{crepe} J. W. Kim, J. Salamon, P. Li, and J. P. Bello. \emph{CREPE: A
Convolutional Representation for Pitch Estimation.} ICASSP, 2018.
arXiv:1802.06182. Cài đặt: torchcrepe,
\url{https://github.com/maxrmorrison/torchcrepe}.

\bibitem{pesto} A. Riou, S. Lattner, G. Hadjeres, and G. Peeters. \emph{PESTO:
Pitch Estimation with Self-supervised Transposition-equivariant Objective.}
International Society for Music Information Retrieval Conference (ISMIR), 2023.
\url{https://github.com/SonyCSLParis/pesto}.

\bibitem{penn} M. Morrison, C. Hsieh, N. Pruyne, and B. Pardo.
\emph{Cross-domain Neural Pitch and Periodicity Estimation.} 2023.
arXiv:2301.12258. Cài đặt: penn,
\url{https://github.com/interactiveaudiolab/penn}.

\bibitem{basicpitch} R. M. Bittner, J. J. Bosch, D. Rubinstein,
G. Meseguer-Brocal, and S. Ewert. \emph{A Lightweight Instrument-Agnostic Model
for Polyphonic Note Transcription and Multipitch Estimation.} ICASSP, 2022.
\url{https://github.com/spotify/basic-pitch}.

\bibitem{krumhansl} C. L. Krumhansl. \emph{Cognitive Foundations of Musical
Pitch.} Oxford University Press, 1990. (Các hồ sơ âm thử giọng Krumhansl-Kessler /
Krumhansl-Schmuckler.)

\bibitem{librosa} B. McFee, C. Raffel, D. Liang, D. P. W. Ellis, M. McVicar,
E. Battenberg, and O. Nieto. \emph{librosa: Audio and Music Signal Analysis in
Python.} Proceedings of the 14th Python in Science Conference (SciPy), 2015.
\url{https://librosa.org/}.

\bibitem{music21} M. S. Cuthbert and C. Ariza. \emph{music21: A Toolkit for
Computer-Aided Musicology and Symbolic Music Data.} ISMIR, 2010.
\url{https://web.mit.edu/music21/}.

\bibitem{verovio} L. Pugin, R. Zitellini, and P. Roland. \emph{Verovio: A
library for Engraving MEI Music Notation into SVG.} ISMIR, 2014.
\url{https://www.verovio.org/}.

\bibitem{mireval} C. Raffel, B. McFee, E. J. Humphrey, J. Salamon, O. Nieto,
D. Liang, and D. P. W. Ellis. \emph{mir\_eval: A Transparent Implementation of
Common MIR Metrics.} ISMIR, 2014. \url{https://github.com/craffel/mir_eval}.

\bibitem{hollien} Hollien, Dew, and Philips. \emph{Phonational Frequency Ranges
of Adults.} Journal of Speech and Hearing Research, 14:755, 1971. (Chuẩn quần thể
cho bộ tìm quãng giọng: trung bình 38 nửa cung cho nam và 37 cho nữ; độ lệch chuẩn
dùng ở đây là một giả định được nêu rõ, không lấy từ nghiên cứu này.)

\end{thebibliography}

% =========================================================================
% PHU LUC A - MINH HOA HE THONG (ANH CHUP MAN HINH)
% =========================================================================
\newpage
\appendix
\addcontentsline{toc}{section}{Phụ lục A: Minh họa hệ thống}
\section*{Phụ lục A: Minh họa hệ thống}
\renewcommand{\thesubsection}{A.\arabic{subsection}}
\setcounter{subsection}{0}

Phụ lục này trình bày ứng dụng web đang chạy từ góc nhìn người dùng. Mỗi tiểu mục
ứng với một thẻ. Các ảnh chụp được lấy trên một phiên bản cục bộ của ngăn xếp; các
chi tiết hiển thị nhỏ tùy thuộc trình duyệt và hệ điều hành của người dùng.

\subsection{Hub}

Hub là thẻ trang chủ: một bảng luyện tập chỉ đọc, gộp lịch sử mà các thẻ khác ghi
lại cục bộ thành các đường xu hướng ngoại tuyến và một báo cáo tiến bộ gọn. Thẻ Coach
của nó thêm nút tùy chọn ``Nhận một kế hoạch luyện tập''.

\begin{figure}[H]
\centering
\reportfig{0.85}{screenshot-hub.png}{Thẻ Hub: các xu hướng tiến bộ ngoại tuyến gộp
từ lịch sử luyện tập cục bộ, và thẻ tùy chọn ``Nhận một kế hoạch luyện tập''.}
\caption{Bảng luyện tập Hub.}
\label{fig:scr-hub}
\end{figure}

\subsection{Bộ ký âm: chế độ Auto}

Thẻ Bộ ký âm ghi âm hoặc tải lên một đoạn ngân nga và trả về bản nháp tự động: bản
nhạc được khắc, giọng phát hiện được, và các hợp âm gợi ý, kèm nút tải MIDI và
MusicXML. Các hợp âm có thể được làm lại tại chỗ: nhấp một hợp âm đưa ra các lựa
chọn ngoại tuyến, và một hàng ``Reharmonise'' hỏi mô hình ngôn ngữ tùy chọn một tiến
trình mới theo một phong cách yêu cầu.

\begin{figure}[H]
\centering
\reportfig{0.85}{screenshot-transcriber.png}{Thẻ Bộ ký âm ở chế độ Auto: bản nhạc
khắc phía máy chủ, giọng phát hiện được, và các nút xuất.}
\caption{Bộ ký âm: chế độ Auto.}
\label{fig:scr-transcriber}
\end{figure}

\subsection{Bộ ký âm: chế độ Manual}

Chế độ Manual là trình soạn thảo khuông nhạc trong trình duyệt để sửa bản nháp. Dải
cao độ tham chiếu hiển thị đường cong ngân nga so với các nốt đã chọn, và thanh công
cụ chỉnh sửa cung cấp đổi cao độ, trường độ, gộp, tách, xóa và chèn kèm hoàn tác và
làm lại. Một ô ``Ask'' cũng nhận các chỉnh sửa bằng ngôn ngữ tự nhiên, xử lý bằng
một văn phạm ngoại tuyến cho các cách diễn đạt phổ biến và mặt khác bằng mô hình
ngôn ngữ tùy chọn.

\begin{figure}[H]
\centering
\reportfig{0.85}{screenshot-manual.png}{Thẻ Bộ ký âm ở chế độ Manual: khuông nhạc
chỉnh sửa được, dải cao độ tham chiếu, và thanh công cụ chỉnh sửa.}
\caption{Bộ ký âm: chế độ Manual.}
\label{fig:scr-manual}
\end{figure}

\subsection{Bộ tìm cao độ}

Thẻ Bộ tìm cao độ phân tích một đoạn tải lên để tìm giọng, nhịp độ và một mã Camelot,
dùng phương pháp dựa trên chroma hoạt động được trên nhạc đa âm.

\begin{figure}[H]
\centering
\reportfig{0.85}{screenshot-pitchfinder.png}{Thẻ Bộ tìm cao độ: giọng phát hiện
được, nhịp độ, mã Camelot, và các thống kê kỹ thuật cho một đoạn tải lên.}
\caption{Bộ tìm cao độ.}
\label{fig:scr-pitchfinder}
\end{figure}

\subsection{Thời gian thực và luyện giọng}

Thẻ Realtime là một bộ giám sát cao độ phía máy khách và bộ lên dây đàn guitar với
các bài luyện giọng: một làn nốt mục tiêu, một trò chơi ghép nốt, một bộ luyện thang
âm, phân tích vibrato, một bộ tìm quãng giọng, và một bảng tiến bộ theo phiên.

\begin{figure}[H]
\centering
\reportfig{0.85}{screenshot-realtime.png}{Thẻ Realtime: bộ giám sát giọng trực tiếp
với bộ chọn chế độ luyện tập (Free / Match game / Scale trainer / Range), đồng hồ
lên dây, chỉ số độ ổn định, và đồ thị cao độ trực tiếp.}
\caption{Bộ giám sát giọng thời gian thực và luyện giọng.}
\label{fig:scr-realtime}
\end{figure}

\subsection{Bộ dịch giọng}

Thẻ Bộ dịch giọng chuyển một bản nhạc sang giọng mới. Chế độ tệp dịch một tệp MIDI
hoặc MusicXML tải lên ở phía máy chủ và cung cấp các giọng tương thích Camelot dạng
phím dịch; chế độ ngân nga dịch bản ký âm gần nhất ở phía máy khách.

\begin{figure}[H]
\centering
\reportfig{0.85}{screenshot-transposer.png}{Thẻ Bộ dịch giọng ở chế độ ngân nga: giai
điệu vừa ngân nga được dịch phía máy khách, kèm nút dịch, bộ chọn giọng đích, và các
hợp âm đã dịch. Chức năng hợp âm (số La Mã) được bảo toàn, nên chỉ tên hợp âm thay
đổi.}
\caption{Bộ dịch giọng: chế độ ngân nga.}
\label{fig:scr-transposer}
\end{figure}

\subsection{Sing-Along và gợi ý luyện tập bằng AI}

Thẻ Sing-Along phát một bản nhạc tải lên làm hướng dẫn, theo dõi giọng trực tiếp, và
chấm điểm từng nốt. Sau một lượt hát, nó hiển thị một phân tích phần trình bày, xác
định, ngoại tuyến, và, khi được yêu cầu, gửi một bản tóm tắt số tới một mô hình ngôn
ngữ bên ngoài để lấy gợi ý luyện tập bằng ngôn ngữ (một trong các tính năng tùy chọn,
không cục bộ mô tả ở Mục~\ref{sec:coaching}).

\begin{figure}[H]
\centering
\reportfig{0.85}{screenshot-singalong.png}{Thẻ Sing-Along sau một lượt hát: tổng
quan chấm điểm từng nốt, phân tích xác định (độ lệch cao độ, độ trôi, độ chính xác
bước nhảy so với bước liền, quãng giọng yếu nhất, các nốt tệ nhất), và bảng tùy chọn
``Nhận gợi ý luyện tập'' kèm bộ chọn ngôn ngữ.}
\caption{Sing-Along và gợi ý luyện tập bằng AI.}
\label{fig:scr-singalong}
\end{figure}

\subsection{Bộ luyện tai}

Thẻ Bộ luyện tai phát một nhóm nốt và yêu cầu người dùng nhận diện nó bằng cách chạm
một nút trắc nghiệm, với các chế độ con cho quãng, chất lượng hợp âm, thang âm và
điệu thức, và bậc thang âm, ở ba mức độ khó. Nó hoàn toàn ở phía máy khách, không
micro và không máy chủ.

\begin{figure}[H]
\centering
\reportfig{0.85}{screenshot-eartrainer.png}{Thẻ Bộ luyện tai: bộ chọn chế độ con và
độ khó, nút phát, các nút đáp án trắc nghiệm, và điểm số, chuỗi đúng, cùng lịch sử
độ chính xác theo từng chế độ đang chạy.}
\caption{Bộ luyện tai.}
\label{fig:scr-eartrainer}
\end{figure}

\end{document}
